 \documentclass[aps,prx,reprint,superscriptaddress,nofootinbib,longbibliography]{revtex4-1}

\usepackage{amsmath,amssymb, amsthm,amstext}
%\usepackage{booktabs}
%\usepackage{caption}
\usepackage{natbib}
\usepackage{graphicx}
\usepackage{color}
\usepackage{array, enumerate}
\usepackage{bm}
\usepackage{multirow}
\usepackage[breaklinks,colorlinks,citecolor=blue]{hyperref}
\usepackage{braket}
\usepackage{txfonts}
\usepackage{physics}
\usepackage{placeins}
\usepackage[normalem]{ulem}


\newcommand{\todo}[1]{\textcolor{red}{#1}}

\def\nn{\nonumber}
\def\be{\begin{equation}}
\def\ee{\end{equation}}


\def\apjl{Astroph.J.Lett.}
\def\apjs{Astroph.J.S.}
\def\mnras{MNRAS}
\def\araa{ARA\&A}
\def\aap{Astronomy\&Astrophysics}
\def\baas{Bulletin of the AAS}


\begin{document}

%\title{Secular evolution of extreme mass ratio inspirals in a wet environment: \\ A relativistic treatment of collisions with disks}
\title{Captured are circularized:  A relativistic \\ treatment  of extreme mass ratio inspirals crossing accretion disks}
\author{Yuhe Zeng}
\email{zengyuhe@sjtu.edu.cn}
\affiliation{Tsung-Dao Lee Institute, Shanghai Jiao-Tong University, Shanghai, 520 Shengrong Road, 201210, People’s Republic of China}
\affiliation{School of Physics \& Astronomy, Shanghai Jiao-Tong University, Shanghai, 800 Dongchuan Road, 200240, People’s Republic of China}


\author{Zhen Pan}
\email{zhpan@sjtu.edu.cn}
\affiliation{Tsung-Dao Lee Institute, Shanghai Jiao-Tong University, Shanghai, 520 Shengrong Road, 201210, People’s Republic of China}
\affiliation{School of Physics \& Astronomy, Shanghai Jiao-Tong University, Shanghai, 800 Dongchuan Road, 200240, People’s Republic of China}


\begin{abstract}
 A small body orbiting around an accreting massive object and periodically crossing its accretion disk is a common configuration in astrophysics. In this work, we investigate the secular evolution of extreme mass-ratio inspirals (EMRIs), in which a stellar-mass object (SMO), e.g., a star or a stellar-mass black hole (sBH), collides with the accretion disk of a central supermassive black hole (SMBH), within a fully relativistic framework. We find (1) the disk always tends to align the SMO no matter what the initial orbital inclination $\iota$ relative to the disk is, (2) the final orbital eccentricity
 of the SMO captured by the disk is always low though the orbital eccentricity may temporarily grow when the orbital inclination $\iota$ is large and the SMO is an sBH,
 and (3) via collisions with the accretion disk only, only a small fraction of sBHs that are initially close to the SMBH and close to the disk can be captured by the disk within typical disk lifetime of active galactic nuclei. Two-body scatterings between SMOs in the nuclear stellar cluster play an essential role in randomly kicking sBHs towards the disk and boosting the capture rate.
\end{abstract}
\date{\today}

\maketitle

\section{Introduction}
%{\bf 1st paragraph: briefly summarize astrophysical systems consisting of a small body orbiting around an accreting massive body, and previous studies on the secular evolution }
%{\bf 2nd paragraph: turn to EMRI discussion, brief background review, why a relativistic treatment is needed. QPEs clearly belong to the topic of this work and is a motivation of this work}
%{\bf Be more careful with your references: e.g., Dong Lai's paper I sent you earlier should be cited in the first paragraph, where they studied planet disk collisions. More relevant than embedded systems.}
Astrophysical systems in which a small body orbits around an accreting massive central object appear in a wide range of environments, e.g.,  (i) planetesimals interacting with circumstellar material around  white dwarfs (WDs) or young planets embedded in protoplanetary disks, where gas drag or disk–planet torques drive orbital migration and damp eccentricity and inclination~\cite{OConnor2020,annurev:/content/journals/10.1146/annurev-astro-081811-125523,refId1}; (ii) compact stellar-mass objects (SMOs) or stars interacting with the dense gas of active galactic nuclei (AGN) disks, which can trap, align and migrate small objects and facilitate hierarchical mergers~\cite{Yang_2019,Peng_2025,10.1111/j.1745-3933.2011.01132.x,10.1111/j.1365-2966.2006.11155.x,10.1093/mnras/stw2260,Tagawa_2020}; (iii) SMOs around supermassive black holes (SMBHs) interacting with transient accretion disks formed in tidal disruption events (TDEs) and producing X-ray  quasi-periodic eruptions (QPEs)~\cite{Xian_2021,Linial_2023,refId2,10.1093/mnras/stad2616,PhysRevD.109.103031,PhysRevD.110.083019,Zhou_2025}.
%{\bf [Xian+2021,Linial+2023,Franchini+2023,Tagawa+2023,Zhou+2024-2025]}.
Together, these systems provide a unified physical setting in which an orbiting secondary interacts with the central massive object itself or with its ambient environment, leading to secular modifications of the orbit and distinctive radiative signatures. A substantial body of work has focused on how such interactions drive the secular evolution of orbiting secondaries. Depending on the environment, long-term orbital changes may arise from gravitational radiation, dissipative gas torques, dynamical friction~\cite{refId0,spieksma2025gripdiskdraggingcompanion}, or impulsive perturbations during disk crossings~\cite{OConnor2020,Wang2024,spieksma2025gripdiskdraggingcompanion,Rowan_2025,10.1093/mnras/staa3004,10.1093/mnras/stad1295,Secunda_2021,MacLeod:2019jxd}. Over timescales much longer than the orbital period, the small effects produced by these forces accumulate into a net secular drift of the orbital elements, corresponding to the usual adiabatic evolution of the system.
%Over timescales much longer than the orbital period, these effects accumulate coherently and reshape the orbital elements in a systematic, adiabatic manner. {\bf This last sentence is obviously written by some large language model. Sometime it does not understand what it is writing. You need to be responsible for what it is writing for your paper: what's the physical meaning of coherently and systematic here ?}



 One class of such astrophysical systems, extreme mass-ratio inspiral (EMRI), typically consisting of a central SMBH and an orbiting SMO, which can be a stellar companion or a compact remnant such as a stellar-mass black hole (sBH), satisfying a mass ratio $q\equiv m_{\rm SMO}/M_{\rm SMBH}\lesssim 10^{-4}$, is among the primary gravitational-wave sources targeted by the future space-based missions LISA~\cite{baker2019laserinterferometerspaceantenna}, TianQin~\cite{10.1093/ptep/ptaa114} and Taiji~\cite{10.1093/ptep/ptaa083}. Their long-lived, relativistic dynamics encode detailed information about the strong-field spacetime and thereby enable precision tests of general relativity (GR) and measurements of SMBH properties~\cite{AmaroSeoane2018,Barack:2018yvs,Gair:2012nm}. Although EMRIs are often treated as vacuum systems, realistic AGNs contain dense stellar clusters and accretion disks, within which compact objects may reside or migrate~\cite{PhysRevD.75.024034,10.1111/j.1365-2966.2006.11155.x,PhysRevD.103.103018,PhysRevD.104.063007,Pan_2021,PhysRevD.105.083005,10.1093/mnras/stad749}.
 %[ {\bf (Sigl et al. 2007; Levin 2007; Pan \&657Yang 2021a; Pan et al. 2021; Pan \& Yang 2021b; Pan et al.6582022; Derdzinski \& Mayer 2023) these are more relevant refs already cited in the secular evolution of QPEs paper}].
 In such environments, SMOs in EMRIs may experience non-negligible matter-induced forces, including aerodynamic drag, hydrodynamical torques, and gas dynamical friction, all of which may significantly modify their orbital evolution relative to vacuum inspirals. Such systems, often termed ``wet EMRIs"~\cite{PhysRevD.104.063007}, can deviate in both secular orbital-parameter evolution and inspiral timescales, affecting not only the astrophysical interpretation of their dynamics but also the accuracy of waveform templates required for LISA data analysis~\cite{PhysRevD.84.024032,10.1093/mnras/staa3976,10.1093/mnras/stz1026}.
 %~\cite{PhysRevD.84.024032,10.1093/mnras/staa3976,10.1093/mnras/stz1026,10.1093/mnras/stae1764}[{\bf Check these refs. Some of them are irrelevant.}].

Formation of wet EMRIs in AGN disks by capturing SMOs on misaligned orbits has been investigated in detail \cite{PhysRevD.105.083005}. 
Two-body scatterings randomly kick SMOs onto low-inclination orbits, and  density waves excited by low-inclination SMOs along with dynamical friction efficiently capture the SMOs onto the AGN disk. Density waves drive the SMOs to migrate inward as well as damping their orbital eccentricities~\cite{Tanaka_2002,Tanaka_2004}. As a result, wet EMRIs in the LISA sensitivity band are expected to be of low eccentricity. On the other hand, there are also recent studies~\cite{10.1093/mnras/stae1239,PhysRevD.111.084006} of eccentric wet EMRIs and their scientific potentials motivated by previous studies showing that SMO-disk collisions may increase the orbital eccentricity in some parameter space~\cite{Secunda_2021,Wang2024}.
% [\textcolor{red}{No observation for eccentricity excitation???}].
 
% In the star-disk collision scenario of a wet EMRI as depicted in Fig.~\ref{fig:EMRI scheme}, the SMO with an inclined orbit relative to the disk plane periodically crosses the accretion disk surrounding the SMBH, which arises naturally in dense stellar clusters where compact objects gradually migrate towards the SMBH~\cite{Amaro-Seoane_2007}, and can cumulatively modify the SMO's orbital energy (corresponding with semi-major axis $a$), eccentricity $e$, and inclination angle $\iota$ over many periods~\cite{10.1093/mnras/stad1016,Wang2024,spieksma2025gripdiskdraggingcompanion,10.1093/mnras/staa3004,10.1093/mnras/stad1295}. The efficiency of these interactions depends sensitively on the properties of the accretion disk, primarily its density profile~\cite{1973A&A....24..337S,Thompson_2005}.
% %Consequently, such EMRIs provide a natural laboratory to study the interplay between relativistic orbital dynamics and gas-induced perturbations, bridging the gap between idealized vacuum EMRI models and more realistic astrophysical systems.
% Previous studies of star-disk interactions in EMRIs have primarily focused on simplified models in which the SMO interacts with the accretion disk through phenomenological drag or torque prescriptions, and based on the simulation approaches assuming Newtonian or post-Newtonian approximations and neglect fully relativistic effects, which can become significant in the strong-field region near the SMBH~\cite{Wang2024,spieksma2025gripdiskdraggingcompanion,10.1093/mnras/staa3004,10.1093/mnras/stad1295,Kocsis:2011dr,PhysRevLett.107.171103}. Once the orbital radius approaches the innermost few tens of gravitational radii, a fully relativistic framework will be gradually essential for reliably characterizing the secular evolution induced by star–disk collisions in EMRIs, ensuring that the underlying orbital dynamics can be treated well in the regime most relevant for LISA-band gravitational-wave signals, and also enabling more reliable estimates of the characteristic timescales for disk-assisted orbital capture or shrinkage and of the accompanying evolution of the orbital eccentricity. For QPEs modeled in the context of star-disk collision within EMRI systems, the timing, luminosity, and energy spectrum of the flares may be highly sensitive to the relativistic orbital dynamics. This strongly motivates the use of a fully relativistic framework to model the secular evolution under star–disk collisions, in order to produce accurate predictions for flare recurrence, energetics, and the evolution of orbital parameters-predictions that can be directly confronted with observed QPE light curves.\par
In this work, we investigate the orbital evolution induced by SMO-disk collisions  within a fully relativistic framework. Specifically, we model the motion of the SMO as a forced geodesic in Schwarzschild spacetime and consider its repeated crossings of a geometrically thin accretion disk located on the equatorial plane (see Fig.~\ref{fig:EMRI scheme}). To quantify the secular evolution of the orbit, we derive the equations governing the variations of the principal orbital elements $a$, $e$, and $\iota$, directly with respect to the proper time $\tau$. By applying an adiabatic approximation procedure over the fast orbital phases, we extract the net secular evolution driven by SMO-disk collisions and determine the associated characteristic timescales for orbital shrinkage, inclination damping, and also the tendency of eccentricity evolution.\par

This paper is organized as follows. 
In Sec.~\ref{sec:basic framework}, we introduce the evolution equations for the orbital elements, the adiabatic orbit-averaging scheme used to obtain their secular behavior, and also summarize the two force models considered in this work.
In Sec.~\ref{sec:results}, we present the secular evolution of the orbital parameters under the two force models and derive the scaling relations that govern the characteristic timescales, and compare our findings with previous studies, extend the timescale estimates using the adopted disk model and scaling relations.
We conclude this work with  Sec.~\ref{sec:conclusion}. In Appendices~\ref{appendix:form of acceleration}, \ref{appendix:coefficients}, and \ref{appendix:average Integral2}, we provide details of the key formulae used in this work. Throughout this work, we use geometrized units with $G=c=1$.

\begin{figure*}[t]
    \centering
    \includegraphics[scale=0.45]{small_figures/Scheme.jpg}
    \caption{Schematic illustration of the wet EMRIs with SMO-disk collision scenario. The relations between the labeled parameters are $z_{1}\equiv\cos\theta_{\rm min}=\sin\iota$, where $\Vec{L}_{\rm disk}$ is the disk angular momentum, and $\Vec{L}_{\rm orbit}$ is the orbital angular momentum of the SMO, the angle between these two vectors is equivalent to the orbital inclination, $\cos^{-1}(\hat{L}_{\rm disk}\vdot\hat{L}_{\rm orbit})=\iota$. $0<\iota<90^{\circ}$ corresponds to the prograde cases, while $90^{\circ}<\iota<180^{\circ}$ corresponds to the retrograde cases. In both cases $\iota$ will decrease, which will be verified in this work.}
    \label{fig:EMRI scheme}
\end{figure*}

\section{Basic framework}
\label{sec:basic framework}
\subsection{Forced EMRI motion around a SMBH}\label{sec:forced motion}
Within the GR framework, the motion of a freely falling test particle is simply described by a geodesic of the background spacetime. 
%, determined entirely by the metric of the central massive object. 
However,  the geodesic motion is perturbed by external forces in non-vacuum environments and the equation of motion (E.O.M.) is written as 
%In such cases, the particle follows a forced geodesic, in which the external forces introduce a non-vanishing four-acceleration and drive secular changes in the orbital parameters. To compute these dynamical effects consistently, one must begin with the general form of the equations of motion for a particle subject to an arbitrary four-acceleration, from which the explicit equation is obtained as
\begin{equation}
    \dv[2]{x^{\mu}}{\tau}+\Gamma^{\mu}_{\alpha\beta}\dv{x^{\alpha}}{\tau}\dv{x^{\beta}}{\tau}=f^{\mu},
    \label{eq: forced geodesic}
\end{equation}
where $x^{\mu}$ is the four-coordinate of the SMO around the central SMBH, and $f^{\mu},\ \tau$ denote the four-acceleration and proper time, respectively. In the vacuum environment, the external force vanishes,  $f^{\mu}=0$. 
%However, the interaction must be considered in the presence of surrounding matter, such as accretion-disk material or a spherically distributed dark-matter (DM) component. 
The form of $f^{\mu}$ in the reference of the zero-angular-momentum observers (ZAMOs) has been given in Ref.~\cite{PhysRevD.83.044037}, $f^{\mu}$ was assumed to depend linearly on the spatially projected component of the relative four–velocity $u_{\perp}^{\mu}$, and in order to enforce the orthogonality condition, additional correction terms associated with $u_{\rm ZAMO}^{\mu}$ were also included, thus,
\begin{equation}
    f^{\mu}=-\gamma u^{\mu}_{\perp}+\kappa u^{\mu}_{\rm ZAMO},
\end{equation}
where $u^{\mu}_{\perp}=u^{\mu}+\Gamma u^{\mu}_{\rm ZAMO}$ with $\Gamma\equiv u^{\mu}_{\rm ZAMO}u_{\mu}$, and with the constraints of $u^{\mu}u_{\mu}=-1$ and equivalently, $f^{\mu}u_{\mu}=0$, there is $\kappa=\gamma(\Gamma-1/\Gamma)$, then
\begin{equation}
    \begin{aligned}
        f^{\mu}&=-\gamma u_{\perp}^{\mu}+\gamma\left(\Gamma-\frac{1}{\Gamma}\right)u^{\mu}_{\rm ZAMO}\\
        &=-\gamma\left(u^{\mu}+\frac{u^{\mu}_{\rm ZAMO}}{u_{{\rm ZAMO}\ t} u^{t}}\right),
    \end{aligned}
    \label{eq: force form}
\end{equation}
with the four velocity of the ZAMOs satisfying $u_{\rm ZAMO\ \phi}=0$. Here $\gamma$ denotes the interaction (damping) coefficient that encodes the strength of the force exerted on the SMO. When combined with the relative velocity between the SMO and the ZAMO frame, it determines the individual components of the resulting four-acceleration.

%{\bf Rewrite this whole paragraph. Deepseek is useful.}
In the case of an SMO interacting with an accretion disk, the disk is assumed to follow nearly Keplerian rotation around the central SMBH and to be geometrically thin. The force exerted on the SMO therefore depends not only on the local disk properties but also on the relative velocity between the SMO and the disk material. In this work, we incorporate both the spatial dependence of the interaction coefficient and the actual relative velocity determined by the Keplerian motion of the disk material within the thin-disk approximation. Starting from the density profile of the disk, we adopt an axisymmetric configuration in which the thin-disk density distribution is expressed as a function of the radial coordinate $r$ and the polar angle $\theta$, assuming a Gaussian density profile in the vertical direction:
\begin{equation}
    \rho_{\rm gas}(r,\ \theta)= \rho_{\rm g} \exp\left[-\frac{(\theta-\pi/2)^2}{2(H_{\rm disk}/r)^2}\right],
    \label{eq: rho gas disk}
\end{equation}
where $\rho_{\rm g}$ is the reference density of the disk and $H_{\rm disk}$ denotes the disk scale height to the midplane. This distribution effectively confines the disk material to the vicinity of the equatorial plane when $H_{\rm disk}\ll r$. Consequently, in our calculations the SMO interacts significantly with the gas only when it passes near the midplane. Moreover, since the present analysis is restricted to a finite radial interval and aims to isolate the dynamical impact of disk crossings rather than to model the global disk structure, we omit any additional radial dependence beyond the exponential angular profile introduced above.
Under the thin-disk approximation, the disk matter is confined to $\theta=\pi/2$, and the corresponding four-velocity components are
\begin{equation}
\begin{aligned}
     &u_{\rm disk}^{r}=u_{\rm disk}^{\theta}=0, \\ 
    &g_{\mu\nu}u_{\rm disk}^{\mu}u^{\nu}_{\rm disk}=-1, \\
    &u_{\rm disk}^{\phi}=\Omega_{K}u_{\rm disk}^{t},
\end{aligned}
\end{equation}
where 
\begin{equation}
    \Omega_{K}^{\pm}=\pm\frac{\sqrt{M_{\bullet}}}{r^{3/2}\pm a_{\bullet}\sqrt{M_{\bullet}}}
\end{equation}
denotes the Keplerian angular velocity in Kerr spacetime with $a_{\bullet}$ the spin parameter of the central SMBH.
The ``+" branch corresponds to prograde rotation, i.e., the angle between the disk angular momentum and the SMBH spin is less than $90^{\circ}$, whereas the ``-" branch corresponds to retrograde rotation. In the present work, which is carried out in Schwarzschild spacetime, the sign is determined by the vertical projection of the orbital angular momentum of the SMO onto the disk angular momentum. By substituting $u^{\mu}_{\rm ZAMO}$ in Eq.~\eqref{eq: force form} with $u^{\mu}_{\rm disk}$, the four-acceleration for the SMO-disk collision scenario is derived, with the $r,\ \theta,\ \phi$ components presented in Appendix~\ref{appendix:form of acceleration}.
%and “+” denotes alignment between the disk angular momentum and the SMBH's spin angular momentum, while in this work, the vertical projection of the SMO’s orbital angular momentum, whereas “–” denotes the opposite orientation. Once the four-velocity of the disk material is determined, the relative velocity between the SMO and the disk can be computed, which in turn yields the four-acceleration acting on the SMO. Detailed forms of $f^{r},\ f^{\theta},\ f^{\phi}$ are provided in Appendix~\ref{appendix:form of acceleration}.

\subsection{Damping coefficient $\gamma$}\label{ref:disk coefficient}
%In Sec.~\ref{sec:forced motion}, the thin disk matter distribution has been shown with exponential form, with the disk drag force coefficient $\gamma$ to be estimated. Refs.~\cite{PhysRevD.109.103031,PhysRevD.110.083019,Zhou_2025} calculated the orbital energy dissipation of the SMO, which is either a star or a small black hole (sBH), after each collision with the disk. In the case that the SMO is a star, Ref.~\cite{OConnor2020} has given the acceleration caused by the drag force in the framework of Newtonian dynamics as
In this subsection, we compute the damping coefficient $\gamma$ in Eq.~\eqref{eq: force form} for both star-disk collisions and sBH-disk collisions. The coefficient $\gamma$ depends on several factors, including the interaction type, the size of the SMO, the gas density and scale height of the disk.\par
According to the thin-disk density profile in Eq.~\eqref{eq: rho gas disk}, the effects of SMO–disk collisions are strongly localized near the equatorial plane, $\theta=\pi/2$. In this region, where $\abs{\theta-\pi/2}\ll 1$, one may approximate $r^2(\theta-\pi/2)^2/H_{\rm disk}^2\approx r^2\cos^2\theta/H_{\rm disk}^2$.
In terms of disk surface density $\Sigma_{\rm g}$ and disk height to the midplane $H_{\rm disk}$, the reference gas density $\rho_{\rm g}$ in Eq.~\eqref{eq: rho gas disk} is estimated by treating the thin disk as having an effective vertical thickness $2H_{\rm disk}$ and a uniform density within this narrow region as  $\rho_{\rm g}\approx \Sigma_{\rm g}/2 H_{\rm disk}$. The damping coefficient is then formulated as 
\begin{equation}
    \gamma(r,\ \theta)=\gamma_{\rm g}\exp\left[-\frac{(\theta-\pi/2)^2}{2(H_{\rm disk}/r)^2}\right],
\end{equation}
where $\gamma_{\rm g}$ is the reference damping coefficient apart from the spatial-distribution dependence.\par
%{\bf A problem : $\Sigma \neq \int \rho r d\theta$ ? \textcolor{red}{For this problem, it is a simplified model to mark the region of interaction, but not the real disk model. In numerical calculations, when $\theta-\pi/2$ is not very small, the exponential part will be calculated as $0$, so it is just useful for thin-disk assumption.}}{\bf Check whether this is a convention used in literature. Don't invent a convention especially when there is some intrinsic inconsistency.}
For star-disk collisions, the perturbation force is dominated by aero-drag force~\cite{Wang2024}
\begin{equation}
    \boldsymbol{F}_{\rm aero} = -\pi R_{\star}^2\rho_{\rm g}v_{\rm rel}\boldsymbol{v}_{\rm rel},
    \label{eq:aero drag}
\end{equation}
where $\boldsymbol{v}_{\rm rel}\equiv \boldsymbol{v}-\boldsymbol{v}_{\rm g}$ denotes the relative velocity between the star and the local gas in the disk,  $v_{\rm rel}\equiv\abs{\boldsymbol{v}_{\rm rel}}$, $R_{\star}$ is the stellar radius and $\rho_{\rm g}$ is the gas density of the disk. Consequently, the reference damping coefficient $\gamma_{\rm g}$ can be written in the numerical form as 
\begin{equation}
\begin{aligned}
    \gamma_{\rm g} &=  \pi R_{\star}^2\rho_{\rm g} v_{\rm rel}m_{\star}^{-1} \\ 
    &= 2.55\times 10^{-8}\left(\frac{R_{\star}} {R_{\odot}}\right)^{2}\left(\frac{m_{\star}}{M_{\odot}}\right)^{-1}\\
    &\cdot\left( \frac{v_{\rm rel}}{0.1 c}\right)
    \left(\frac{\rho_g}{ 
    10^{5}\ \mathrm{g\cdot cm^{-2}}/3M_\bullet}\right) \left[M_{\bullet}^{-1}\right]\ .
    \label{eq:aero-drag force gamma}
\end{aligned}
\end{equation}
%{\bf Complete the sBH-disk collision part following the star-disk collisions}

%{\bf Work harder on the following paragraphs: make them clear and clean. }
For sBH–disk collisions, the effective radius of an sBH with mass $m_{\bullet}$ is given by its horizon radius $R_{\rm H}=2m_{\bullet}$, which is much smaller than that of a stellar object with the same mass, and the perturbation force is dominated by dynamical friction~\cite{Wang2024,Ostriker_1999}
\begin{equation}
    \bm{F}_{\rm dyn}=-\mathcal{I}\frac{4\pi(Gm_{\bullet})^2\rho_{\rm g}}{v_{\rm rel}^{3}}\bm{v}_{\rm rel},
    \label{eq:Fdyn}
\end{equation}
where the coefficient $\mathcal{I}$ is related with Mach number $\mathcal{M}\equiv v_{\rm rel}/c_{s}$. In this work, we note that $\mathcal{I}\sim \mathcal{O}(1)$ in the relevant parameter regime and mainly rescales the interaction strength. We therefore set $\mathcal{I}=1$, for simplicity, which does not affect the qualitative orbital-evolution behavior discussed below. The reference damping coefficient $\gamma_{\rm g}$ associated with the acceleration dominated by dynamical friction can be written in the numerical form as
\begin{equation}
\begin{aligned}
    \gamma_{\rm g}&\approx \frac{4\pi G^2 m_{\bullet}\rho_{\rm g}}{v_{\rm rel}^{3}}\\
    %&= 4.62\times 10^{-15}\ \left(\frac{m_{\star}}{M_{\odot}}\right)\Sigma_{5}H_{1.5M_{\bullet}}^{-1}v_{0.1\ \rm rel}^{-3}[M_{\bullet}^{-1}].
    &= 4.62\times 10^{-15}\ \left(\frac{m_{\bullet}}{M_{\odot}}\right)\left(\frac{v_{\rm rel}}{0.1\ c}\right)^{-3} \left(\frac{\rho_g}{ 
    10^{5}\ \mathrm{g\cdot cm^{-2}}/3M_\bullet}\right) \left[M_{\bullet}^{-1}\right].
    \label{eq:dynamical friction gamma}
\end{aligned}
\end{equation}


%Ref.~\cite{Wang2024} has also given the similar force form for the star-like SMOs under the subsonic case, where the total drag force is dominated by the aerodynamic drag force. 
% In this work, the orbital radius is typically $r\sim 10^{2}\ M_{\bullet}$, for which the velocity of the disk material is $\bm{v}_{\rm d}\sim\sqrt{GM_{\bullet}/r}\sim 0.1\ c$. For a stellar-SMO, the escape speed evaluated using the solar mass $M_{\odot}$ and radius $R_{\odot}$, is $\bm{v}_{\rm ecs}\sim \sqrt{2GM_{\odot}/R_{\odot}}\sim 10^{-3}\ c$, which is much smaller than the relative speed $\abs{\bm{v}_{\rm rel}}\sim 0.1\ c$, Therefore, for a stellar-SMO, the interaction force may be treated using the form in Eq.~\eqref{eq:aero drag}. {\bf why do we care about escape velocity ?}

% According to Ref.~\cite{PhysRevD.103.103018}, the disk surface density can be estimated with a magnitude of $\Sigma_{\rm g}\sim 10^{5}\ \mathrm{g\cdot cm^{-2}}$. For a stellar-SMO of mass $m_{*}$, the damping coefficient $\gamma$ associated with the acceleration dominated by aero-drag force~\eqref{eq:aero drag} can be written in a numerical form as
% \begin{equation}
%     \gamma\sim \pi R_{*}^2\rho_{\rm g} v_{\rm rel}m_{*}^{-1}\sim 2.55\times 10^{-8}\left(\frac{R_{*}}{R_{\odot}}\right)^{2}\left(\frac{m_{*}}{M_{\odot}}\right)^{-1}\Sigma_{5}H_{1.5M_{\bullet}}^{-1} [M_{\bullet}^{-1}],
%     \label{eq:simplified model}
% \end{equation}
% where $\Sigma_{\rm g}=2\rho_{\rm g}H_{\rm disk}$, $H_{\rm disk}$ is the disk height to the midplane, and $\Sigma_{5}\equiv\Sigma_{\rm g}/(10^{5}\ \mathrm{g\cdot cm^{-2}}),\ H_{1.5M_{\bullet}}\equiv H_{\rm disk}/(1.5\ M_{\bullet})$. Consider the more realistic model, $\gamma$ should be related with the relative speed as
% \begin{equation}
%     \gamma\sim 2.55\times 10^{-8}\left(\frac{R_{*}}{R_{\odot}}\right)^{2}\left(\frac{m_{*}}{M_{\odot}}\right)^{-1}v_{\rm rel\ 0.1}\Sigma_{5}H_{1.5M_{\bullet}}^{-1} [M_{\bullet}^{-1}],
%     \label{eq:relative speed aero drag}
% \end{equation}
% where $v_{0.1\ \rm rel}=v_{\rm rel}/0.1\ c$.\par
% %Ref.~\cite{Wang2024} also presented the dynamical friction with the form
For clarity, the reference damping coefficient $\gamma_{\rm g}$ is expressed as the combination of a constant baseline value $\gamma_{0}$ and a velocity–dependent contribution arising from the relative motion. Specifically, for aero-drag force, $\gamma_{\rm g}=\gamma_{0}\ (v_{\rm rel}/0.1\ c)$, whereas for dynamical friction, $\gamma_{\rm g}=\gamma_{0}\ (v_{\rm rel}/0.1\ c)^{-3}$. The relative velocity is computed from 
% the three spatial components of $u_{\rm rel}^{\mu}$ given in Appendix~\ref{appendix:form of acceleration},  under spherical-coordinate approximation
% \begin{equation}
%     v_{\rm rel}=\sqrt{(u^{r}_{\rm rel})^{2}+r^{2}(u^{\theta}_{\rm rel})^2+r^{2}\sin^{2}\theta(u^{\phi}_{\rm rel})^{2}},
% \end{equation}
% which agrees closely with the fully relativistic value
\begin{equation}
    v_{\rm rel}=\sqrt{1-\frac{1}{(u_{ \mu}u^{\mu}_{\rm disk})^2}} \ .
\end{equation}
%having negligible impact on the long-term orbital evolution considered in this work.
%{\bf very confusing description here. I explained to you earlier how to calculate $v_{\rm vel}$ EXACTLY. Explicitly explain in one sentence how to calculate $v_{\rm rel}$. Actually $v_{\rm vel}$ is different at two collision points for an eccentric orbit. How did this subtle point is taken into account ?}

\subsection{Orbital Evolution: Osculating orbit}\label{sec:parameter evolution}
%{\bf Rewrite this paragraph. Pay attention to logic connections between sentences}
There are effectively seven degrees of freedom in the orbital dynamics of an SMO in curved spacetime: four coordinate components and four velocity components subject to one normalization constraint $u^{\alpha} u_{\alpha}=-1$. To parameterize this motion, we introduce a set of five orbital elements, corresponding to the degrees of freedom beyond the $t$ and $\phi$ coordinates, denoted by the conserved parameter set $\mathbf{I}=\{p,\ e,\ z_{1}\}$, and phases $\psi_{r},\ \psi_{\theta}$. These parameters are related to the radial and polar motions in Boyer–Lindquist coordinates through
\begin{equation}
\begin{aligned}
    &r=\frac{p}{1+e\cos\psi_{r}}\ , \\
    &\cos\theta =z_{1}\cos\psi_{\theta}\ .
    \label{eq:r theta Kepler}
\end{aligned}
\end{equation}
Here, $p$ is the semi-latus rectum, $e$ is the eccentricity, and $z_{1}\equiv \cos\theta_{\rm min}$, the angle $\theta_{\rm min}$ is related to the orbital inclination $\iota$ with respect to the disk plane via $\theta_{\rm min}=\abs{\pi/2 - \iota}$. The range $0<\iota<90^{\circ}$ corresponds to prograde orbits, whereas $90^{\circ}<\iota<180^{\circ}$ corresponds to retrograde orbits.  $\psi_{r}$ and $\psi_{\theta}$ represent the phase variables associated with the radial and polar motions, respectively.\par

For simplicity, we consider a Schwarzschild SMBH of mass $M_{\bullet}$ in this work. The spacetime metric is
\begin{equation}
    ds^2=-\left(1-\frac{2M_{\bullet}}{r}\right)dt^2+\left(1-\frac{2M_{\bullet}}{r}\right)^{-1}dr^2+r^2\left(d\theta^2+\sin^2\theta\ d\phi^2\right) \ .
    \label{eq:Schwarzschild metric}
\end{equation}
%Ref.~\cite{PhysRevD.77.044013} developed a hybrid Schwarzschild–post-Newtonian self-force model within the osculating-orbits framework. 
Owing to the spherical symmetry of the Schwarzschild background, for spherically symmetric perturbations, the polar parameter $z_{1}$ remains constant, and the orbital plane may be chosen as the equatorial plane without loss of generality, leaving only a single radial phase parameter $\psi_{r}$ to characterize the motion. However, in the SMO-disk collision scenario, the breaking of spherical symmetry by the existence of equatorial disk requires an additional polar phase parameter $\psi_{\theta}$ to fully characterize the orbital evolution.\par
%the motion was confined to a fixed orbital plane, so the perturbed trajectory could be described using a radial phase variable together with three geometric orbital elements, the semi-latus rectum $p$, the eccentricity $e$, and the value of $\chi$ at periapsis, $w$.
%which significantly increases the computational complexity, the calculation form has been presented in Ref.~\cite{PhysRevD.83.044037}, with the detailed form shown in Appendix~\ref{appendix:coefficients}. \par
% Ref.~\cite{Fujita_2009} has shown the expressions of 4-velocity components in Kerr spacetime. With the spin value $a_{\bullet}=0$, the simplified form in the Schwarzschild spacetime is

Following Ref.~\cite{Fujita_2009}, the equations of motion of a test particle around a Schwarzschild black hole are written as 
\begin{equation}
\begin{aligned}
    \left(\dv{r}{\tau}\right)^{2}&=E^{2}-\left(1-\frac{2M_{\bullet}}{r}\right)\left(1+\frac{J^{2}}{r^{2}}\cdot\frac{1}{1-z_{1}^{2}}\right)\ , \\
    \left(\dv{\theta}{\tau}\right)^{2}&=\frac{J^{2}}{r^{4}\sin^{2}\theta}\cdot\frac{1}{1-z_{1}^{2}}\cdot(z_{1}^{2}-\cos^{2}\theta)\ , \\ 
     \dv{t}{\tau}&=E\left(1-\frac{2M_{\bullet}}{r}\right)^{-1}\ ,\\
      \dv{\phi}{\tau}&=\frac{J}{r^{2}}\cdot\frac{1}{1-\cos^{2}\theta}\ ,
\end{aligned}
\label{eq:four equations of motion}
\end{equation}
where $E,\ J$ are the conservation quantities, energy and angular momentum respectively. In the following derivations, we set $M_{\bullet}=1$ for simplicity.
%while the Carter constant, which is the third conservation quantity, is actually $\mathcal{C}\equiv z_{1}^{2}J^{2}/(1-z_{1}^{2})$ for the inclined orbits in the Schwarzschild spacetime, and this term appears in the second equation of Eqs.~\eqref{eq:four equations of motion}. {\bf where is Carter constant ? \textcolor{red}{There was a typo, $\mathcal{C}$ should appear in the second equation of Eq.~\eqref{eq:four equations of motion}}.} 

Based on Eqs.~\eqref{eq:four equations of motion}, 
%{\bf the fundamental frequency parts ???} 
the frequencies of the unperturbed geodesic motion for phases $\psi_{r},\ \psi_{\theta}$ can be obtained as
%[{\bf Refs ?}]
\begin{equation}
\begin{aligned}
     \dv{\psi_{r}}{\tau}&:=\omega_{r}=\frac{1}{(\partial r/\partial\psi_{r})}\dv{r}{\tau}=(1+e\cos\psi_{r})^{2}\sqrt{\frac{p-6-2e\cos\psi_{r}}{p^{3}(p-3-e^{2})}}\ , \\
    \dv{\psi_{\theta}}{\tau}&:=\omega_{\theta}=\frac{1}{(\partial\theta/\partial\psi_{\theta})}\dv{\theta}{\tau}=\frac{(1+e\cos\psi_{r})^{2}}{p\sqrt{p-3-e^{2}}}\ .
    \label{eq:dpsidtau and dchidtau}
\end{aligned}
\end{equation}
%{\bf notice that the evolution of two phases are coupled, in the case of no interaction ??? in the case of no external perturbation ?}
Notice that $\omega_{\theta}$ depends solely on $\psi_{r}$. By integrating $\dv{\psi_{\theta}}{\psi_{r}}=\omega_{\theta}/\omega_{r}$, the relation between two phases is found as
\begin{equation}
    \psi_{\theta}(\psi_{r})=2\sqrt{\frac{p}{p-6-2e}}F\left(\frac{\psi_{r}}{2}\ \Bigg|\ \frac{4e}{6+2e-p}\right)+ \psi_{\theta\ \rm ini},
\end{equation}
where the function $F$ is the incomplete elliptic integral of the first kind,
\begin{equation}
    F(\phi\ |\ m)=\int_{0}^{\phi}\frac{\rm d\theta}{\sqrt{1-m\sin^2\theta}}.
\end{equation}
%However, in scenarios where orbital parameters evolve over time, this analytical relation may lead to significant deviations—even under slow variation of parameters, particularly in the case of star-disk collisions.

The forced motion of an SMO can be obtained by directly  evolving the E.O.M. (Eq.~\eqref{eq: forced geodesic}). Then the orbital parameters $p,\ e,\ z_{1}$ can be extracted by inverting Eqs.~\eqref{eq:four equations of motion} \cite{Fujita_2009}. However, as the evolution timescale increases far exceeding the orbital period, the computational cost and the difficulty of precision control for such method rise rapidly, limiting the applicability to long-time evolution problems.
%{\bf Rewrite this paragraph: The first sentence is telling readers you're going to introduce a faster evolution method focusing on secular evolution in this paragraph, but the remaining part of the whole paragraph is irrelevant to this goal.} 

%This also motivates a second method, namely to average over the small oscillations on the orbital-timescale and to focus on the secular evolution on longer timescales, from which one obtains and directly solves evolution equations for the orbital parameters.
This motivates the development of faster evolution methods that average over the rapid oscillations on the orbital timescale and focus on the secular evolution of the orbital parameters, which can be divided into two main steps. First, the E.O.M. is reformulated in terms of orbital elements $\{p,\ e,\ z_{1}, \ \psi_r, \ \psi_\theta\}$ using the method of osculating orbit. Second, applying the adiabatic approximation to the reformulated E.O.M. and extracting secular evolution of orbital elements as in the following Sec.~\ref{sec:secular evolution}.
% First, the orbital parameters are obtained from the directly computed osculating trajectory. Second, based on these results, the adiabatic-average approximation is applied to describe the long-term secular evolution, which will be presented in the following subsection~\ref{sec:secular evolution}.

Following Refs.~\cite{PhysRevD.83.044037,PhysRevD.77.044013}, the osculation condition can be written as  
\begin{equation}
\begin{aligned}
     x^{\mu}(\tau)&=x_{\rm G}^{\mu}[\tau,\ I_{0}(\tau)]\ , \\ 
    \dv{x^{\mu}}{\tau}&=\pdv{x_{\rm G}^{\mu}}{\tau}[\tau,\ I_{0}(\tau)] \ ,
\end{aligned}
\label{eq:x xG}
\end{equation}
where $x^{\mu}_{\rm G}$ is the geodesic with orbital parameter set $I_{0}=\{p,\ e,\ z_{1},\ \psi_{0},\ \chi_{0}\}$. Here $\psi_{0}$ and $\chi_{0}$ are the “initial” values of  $\psi_r$ and $\psi_\theta$, respectively. With definition $\psi_{r}=\psi-\psi_{0}$ one can decompose the phase evolution in the radial direction as a geodesic part $\psi$ governed by Eqs.~\eqref{eq:dpsidtau and dchidtau}, and a perturbation induced shift in the “initial” phase $\psi_{0}$. In the same way, the phase evolution in the polar direction is decomposed as two parts, $\psi_{\theta}=\chi-\chi_{0}$. Defined in this way, parameters
$I_{0}$ remain conserved for unperturbed geodesic motion, however, in the presence of perturbations, it undergoes a slow secular drift superposed with oscillatory variations on the orbital timescale.  

% The explicit expressions for these contributions are given in Appendix~\ref{appendix:coefficients}. At each instant $\tau$, the worldline is treated as a geodesic tangent to the true trajectory, with instantaneous orbital element set $I_{0}(\tau)$. The orbital phases $\psi$ and $\chi$ are then evolved by integrating Eq.~\eqref{eq:dpsidtau and dchidtau}, where the corresponding frequencies $\omega_{r}$ and $\omega_{\theta}$ are evaluated using the instantaneous osculating orbital elements. 
%Since our primary interest lies in the secular evolution of $I^{a}(\tau)$ on timescales much longer than the orbital period, this phase-evolution prescription provides adequate accuracy for the purposes of this work. {\bf No approximation is made here. So there is no accuracy problem in principle right ? Why adequate accuracy for the purposes of this work ? }
%but evolves slowly in the presence of perturbations. 
%The orbital phases $\psi$ and $\chi$ are obtained by integrating Eq.~\eqref{eq:dpsidtau and dchidtau} with the slowly-varying parameter set $I^{a}(\tau)$.
%{\bf They are not really slow-varying parameters: slow secular evolution +  oscillations on orbital timescale. Be more clear on how to evolve  $\psi$ and $\chi$ }
%{\bf Explain why use $\psi_0$ and $\chi_0$, instead of $\psi$ and $\chi$, and how to evolve $\psi$ and $\chi$}

With the osculation condition, one can reformulate the E.O.M. ~\eqref{eq: forced geodesic} in terms of parameters $I_0$.
At each instant $\tau$, the worldline is treated as a geodesic tangent to the true trajectory, with instantaneous orbital element set $I_{0}(\tau)$.
Considering 4-velocity $u^\mu$ along with the 1st osculation condition in  Eq.~\eqref{eq:x xG}, one can find \cite{PhysRevD.77.044013}
\begin{equation}
    u^{\mu}:=\dv{x^{\mu}}{\tau}=\pdv{x^{\mu}_{\rm G}}{\tau}+\pdv{x^{\mu}_{\rm G}}{I_{0}^a}\dv{I_{0}^a}{\tau}=\dv{x^{\mu}_{\rm G}}{\tau},
    \label{eq:osculation 0}
\end{equation}
which in combination with the 2nd osculation condition in  Eq.~\eqref{eq:x xG} immediately leads to 
\begin{equation}
    \pdv{x^{\mu}_{\rm G}}{I_{0}^{a}}\dv{I_{0}^a}{\tau}=0 \ .
    \label{eq:osculating a}
\end{equation}

Considering 4-acceleration $\mathrm{d}u^\mu/\mathrm{d}\tau$ along with  Eqs.~(\ref{eq:osculation 0}, \ref{eq:osculating a}) and 
the geodesic definition of $x^{\mu}_{\rm G}$ 
\begin{equation}
    \pdv[2]{x_{\rm G}^{\mu}}{\tau}+\Gamma^{\mu}_{\alpha\beta}\pdv{x_{\rm G}^{\alpha}}{\tau}\pdv{x_{\rm G}^{\beta}}{\tau}=0 \ ,
\end{equation}
the E.O.M. \eqref{eq: forced geodesic} can be rewritten as 
\begin{equation}
     \dv[2]{x^{\mu}}{\tau}=-\Gamma^\mu_{\alpha\beta} u^\alpha u^\beta + f^\mu = \pdv[2]{x_{\rm G}^{\mu}}{\tau}+f^{\mu}\ .
     \label{eq:dv2xmutaufmu}
\end{equation}

% combining this relation with the forced E.O.M.~\eqref{eq: forced geodesic}, together with Eqs.~\eqref{eq:osculation 0} and~\eqref{eq:osculating a}, the second derivative of $x^{\mu}$ can be written as
% \begin{equation}
%     \dv[2]{x^{\mu}}{\tau}=\pdv[2]{x_{\rm G}^{\mu}}{\tau}+f^{\mu}.
% \end{equation}
On the other hand, the same quantity can also be expressed as
\begin{equation}
    \dv[2]{x^{\mu}}{\tau}=\pdv[2]{x_{\rm G}^{\mu}}{\tau}+\pdv{I_{0}^a}\left(\dv{x^{\mu}_{\rm G}}{\tau}\right)\dv{I_{0}^a}{\tau}\ ,
    \label{eq:dv2xmutau}
\end{equation}
using Eqs.~\eqref{eq:x xG} and ~\eqref{eq:osculating a}. 
With Eq.~\eqref{eq:dv2xmutaufmu} and~\eqref{eq:dv2xmutau},  the forced E.O.M. ~\eqref{eq: forced geodesic} can be written as
%{\bf more details should be added here following  \cite{PhysRevD.77.044013}}
\begin{equation}
    \pdv{I_{0}^a}\left(\dv{x^{\mu}_{\rm G}}{\tau}\right)\dv{I_{0}^a}{\tau}=f^{\mu}.
    \label{eq:osculating b}
\end{equation}

The final step is to isolate individual components $d I_0^a/d\tau$ from Eqs.~\eqref{eq:osculating a} and~\eqref{eq:osculating b}, which contains 
eight equations. By imposing the orthogonality condition $f^{\mu}u_{\mu}=0$, and evolving $t$ and $\phi$ explicitly using the geodesic expressions evaluated along the instantaneous orbit~\cite{PhysRevD.83.044037}, three equation components, i.e., $\mu=t,\ \phi$ components of Eq.~\eqref{eq:osculating a} and $\mu=t$ component of Eq.~\eqref{eq:osculating b}, can be removed. %{\bf 8-1-2-1 = 4, why 5 orbital parameters ?}
As a result, the evolution of the five orbital parameters with respect to $\tau$ can be written generally as
\begin{equation}
    \dv{I_{0}}{\tau}=C^{(I_{0})}_rf^{r}+C^{(I_{0})}_{\theta}f^{\theta}+C^{(I_{0})}_{\phi}f^{\phi},
    \label{eq:parameter evolution ode}
\end{equation}
%{\bf notation $C_I(x)\rightarrow C^{(I)}_x$}
with the detailed expressions of coefficients $C^{(I_{0})}_{r,\theta,\phi}$ shown in Appendix~\ref{appendix:coefficients}.
%By continuously updating the slowly evolving parameters $p,\ e$, and $\iota$, the orbital phases $\psi,\ \chi$ can be then be computed exactly using Eq.~\eqref{eq:dpsidtau} and Eq.~\eqref{eq:dchidtau}. 
%{\bf The following two sentences are hard to understand} In fact, when external perturbations are present, the evolution of the phases, and of the coordinate time $t$, should, in principle, include higher-order perturbative contributions in addition to the unperturbed terms. However, explicit calculations show that, within the context of this work, for such fast-varying parameters, as well as the adiabatic-average method to be discussed later, considering the evolution equations without the perturbative terms, while evaluated with the slowly-varying parameters, is already sufficient.
%the fundamental frequencies associated with the slowly varying parameters are already sufficient.

\subsection{Secular evolution}\label{sec:secular evolution}
%In Refs.~\cite{PhysRevD.83.044037,PhysRevD.77.044013,PhysRevD.77.044012}, the average method for the parameter evolution has been shown, for the cases with one phase. However, when there are multiple phases, each averaged frequency should be decoupled, which is also shown in Ref.~\cite{PhysRevD.83.044037}, as
%{\bf Simplify this paragraph. Make your punch point clear. }
%In the cases that modeling orbital motion under the influence of small perturbing forces, although the cumulative perturbed effects over many orbital periods lead to a slow, secular evolution of the orbital parameters, the motion can be well approximated by a geodesic on short timescales. Thus, the adiabatic-average approximation can be used to technically capitalize on the clear separation between the fast orbital timescale and the slow inspiral timescale. The rapidly oscillating components of the perturbing force, which average over an orbital cycle, have a negligible net effect on the long-term evolution of the orbital parameters. For the star–disk collision problem studied in this work, the SMO interacts with the thin disk only twice during each orbital period, and the influence of each encounter remains within the perturbative regime. Consequently, this method can also be employed to compute the secular evolution.\par
In systems subject to weak perturbations, the orbital motion is well approximated by a geodesic on short timescales, while the perturbations accumulate over many orbital periods and induce a slow secular evolution of the orbital parameters. Such separation between the short orbital period and the long inspiral timescale justifies the use of the adiabatic approximation. For star–disk collisions considered in this work, the SMO interacts with the thin disk twice per orbital period, and each encounter produces a small perturbation, allowing this method to reliably determine the secular evolution.\par

%{\bf Read and rewrite the following two paragraphs. It's very hard to understand them. } 
%In Ref.~\cite{PhysRevD.83.044037}, based on the Mino time $\lambda$~\cite{PhysRevD.67.084027}, the general form of the evolution equation $f$ for the slowly varying parameters $I$ was given as 
%In Ref.~\cite{PhysRevD.83.044037}, the general function of the adiabatic-average method has been shown. {\bf The first sentence is confusing.}
Consider a perturbed system characterized by $N$ phases $\boldsymbol{\psi}=\{\psi_{1},\ \psi_{2},\cdots\psi_{N}\}$ and a set of orbital parameters $\mathbf{I}=\{I_1,\ I_2,\cdots I_{M}\}$ that would remain conserved in the absence of perturbations. For a function $f=f(\boldsymbol{\psi},\ \mathbf{I})$, if the unperturbed frequencies associated with each $\psi_{j}$ depend only on that phase, namely $\omega_{j}=\omega_{j}(\psi_{j},\ \mathbf{I})$, and if each $\psi_{j}$ is assumed to have a period of $2\pi$, then the adiabatic approximation of $f(\boldsymbol{\psi},\ \mathbf{I})$ can be evaluated as~\cite{PhysRevD.83.044037}
\begin{equation}
    \expval{f}_{\mathbf{I}}=\frac{\int_{0}^{2\pi}\mathrm{d}^{N}\boldsymbol{\psi}\prod_{j=1}^{N}\frac{1}{\omega_{j}(\psi_{j},\ I)}f(\boldsymbol{\psi},\ \mathbf{I})}{\prod_{j=1}^{N}\int_{0}^{2\pi}\mathrm{d}\psi_{j}\frac{1}{\omega_{j}(\psi_{j},\ I)}},
    \label{eq: average integral}
\end{equation}
where the numerator on the right-hand side denotes the $N$-dimensional integral over the phase set $\boldsymbol{\psi}$. It has been noted in Ref.~\cite{PhysRevD.83.044037} that for more general systems in which $\omega_{j}=\omega_{j}(\boldsymbol{\psi},\ \mathbf{I})$, the adiabatic approximation may not take the form of Eq.~\eqref{eq: average integral}.
%{\bf If the Kerr problem is NOT investigated, remove the following two senstences.} For the Kerr spacetime, when evolved with the Mino time parameter~\cite{PhysRevD.67.084027}, the radial and polar motions decouple~\cite{Fujita_2009,van_de_Meent_2020}. Consequently, the unperturbed frequencies in Eq.~\eqref{two coordinates r} and Eq.~\eqref{eq:two coordinates theta}, associated with the phases $\psi,\ \chi$, respectively, satisfy the separable form required by Eq.~\eqref{eq:average integral1}.

%where $\psi_{j}$ represents rapidly varying parameters with the fundamental frequency $\omega_{j}$. It should be noted that, in principle, when employing such methods, the $j-$th frequency function must depend solely on its own associated rapidly varying parameter $\psi_{j}$, without any coupling to others. For the Kerr spacetime, the use of Mino time decouples the originally coupled coordinate components $r$ and $\theta$, and thereby decouples their associated phases $\psi$ and $\chi$. Consequently, this method can be employed to obtain the secular evolution of the parameters.\par
%However, this method becomes inapplicable in the full Kerr background due to the coupling between the phases if one attempts to study the parameter evolution directly with respect to the proper time $\tau$ or the phases. 
% In Schwarzschild spacetime, %the adiabatic–average method can be applied directly without introducing the Mino time.
%{\bf Comments in the current form on Refs.~\cite{PhysRevD.77.044013,PhysRevD.77.044012} are distracting and are not helping your work}
%In Refs.~\cite{PhysRevD.77.044013,PhysRevD.77.044012}, the evolution of spherically symmetric orbits in a Schwarzschild background was investigated under the constraint $\theta=\pi/2$. In this case, only the phase associated with the radial coordinate $r$ needs to be considered, and Eq.~\eqref{eq: average integral} can be applied with $N=1$. However, in the star–disk collision scenario, the spherical symmetry is broken, and $\theta$ can no longer be restricted to $\pi/2$.
%Examining the unperturbed frequencies in Eqs.~\eqref{eq:dpsidtau and dchidtau}, it can be found that 
In Schwarzschild spacetime, both frequencies in Eqs.~\eqref{eq:dpsidtau and dchidtau}
solely depend on phase $\psi_{r}$. This corresponds to a special case in which the adiabatic average can still be evaluated in a form analogous to Eq.~\eqref{eq: average integral} as 
%In the Schwarzschild spacetime, the fundamental frequencies governing the evolution of the phases $\psi$ and $\chi$ are both determined solely by the radial phase $\psi$ according to Eq.~\eqref{eq:dpsidtau} and Eq.~\eqref{eq:dchidtau}. Therefore, an alternative form of the averaging formula can be employed to compute the secular evolution with respect to the proper time $\tau$ as
\begin{equation}
    \expval{f}_{\mathbf{I}}=\frac{\int_{0}^{2\pi}d\psi_{\theta}\int_{0}^{2\pi}d\psi_{r}\frac{1}{\omega_{r}(\psi_{r},\ \mathbf{I})}f(\psi_{r},\ \psi_{\theta},\ \mathbf{I})}{2\pi\int_{0}^{2\pi}d\psi_{r}\frac{1}{\omega_{r}(\psi_{r},\ \mathbf{I})}},
    \label{eq:average integral1}
\end{equation}
note that the integration is carried out over the periodic domains of both $\psi_{r}$ and $\psi_{\theta}$, while only the $\omega_{r}$ appears in the denominator of the integrand. The equivalent form adopted in this work is shown in Appendix~\ref{appendix:average Integral2}. The adiabatic E.O.M. is obtained by applying the average above to both sides of Eq.~\eqref{eq:parameter evolution ode}, 
\begin{equation}
    \dv{\expval{I_{0}}}{\tau}=\expval{C^{(I_{0})}_rf^{r}+C^{(I_{0})}_{\theta}f^{\theta}+C^{(I_{0})}_{\phi}f^{\phi}}\ ,
    \label{eq:secular ode}
\end{equation}
where  $\expval{C^{(I_{0})}_if^i}$ on the right hand side can be precomputed for the purpose of numerical efficiency.

%{\bf remove this sentence $\rightarrow$} This eliminates the need to decouple the two phases using the Mino time. {\bf $\leftarrow$ not much relevant to your work} By substituting $f$ with $\mathrm{d}I^{a}/\mathrm{d}\tau$ in Eq.~\eqref{eq:parameter evolution ode}, the secular evolution of the orbital parameters $I^a$ can be obtained. 
%For the spherically symmetric orbits in Schwarzschild background studied in Refs.~\cite{PhysRevD.77.044013,PhysRevD.77.044012}, only the phase associated with the radial coordinate $r$ needs to be considered, and Eq.~\eqref{eq:average integral1} reduces to Eq.~\eqref{eq: average integral} with $N=1$.

\section{Results}
\label{sec:results}
% \subsection{Parameter evolution with the simplified model}
% \label{sec:simplified model}
% To assess the reliability of the adiabatic approximation in the star–disk collision scenario, we first compare the evolution of the three orbital parameters, the semi-major axis $a\equiv p/(1-e^2)$, the eccentricity $e$ and the inclination angle $\iota$, under the simplified model which assumes $\abs{\boldsymbol{v}_{\rm rel}}=0.1\ c$ in Eq.~\eqref{eq:aero drag} for a stellar-SMO system. The trajectories are obtained from (i) the parameter–extraction procedure applied to the evolved E.O.M, as described in Sec.~\ref{sec:parameter evolution}, and (ii) the adiabatic approximation method presented in Sec.~\ref{sec:secular evolution}.\par
% Evolution trajectories of the three parameters over coordinate time $t$ are shown in Fig.~\ref{fig:p 200 e 0.3 z1 40 deg three evo}, demonstrating good agreement between the two methods. The damping coefficient $\gamma$ is set to $2\times 10^{-6}\ M_{\bullet}^{-1}$, and the initial conditions of the orbital parameters are $p_{\rm ini}=200\ M_{\bullet},\ e_{\rm ini}=0.3$ and $\iota_{\rm ini}=50^{\circ}$, while for $\psi_{0}$ and $\chi_{0}$, the secular evolution is negligible in the present calculation, and the initial values are taken as $\psi_{0\ \rm ini}=0$ and $\chi_{0\ \rm ini}=\pi/3$. The disk height to the midplane is set to $H_{\rm disk}=1.5\ M_{\bullet}$ as in Refs.~\cite{PhysRevD.110.083019,Zhou_2025}.\par
% As shown in the top panel of Fig.~\ref{fig:p 200 e 0.3 z1 40 deg three evo}, the semi-major axis $a$ decays over time, reflecting the orbital energy dissipation due to the SMO-disk collision processes. In the middle panel, the eccentricity $e$ also decays with time, indicating that this interaction drives orbital circularization. In the bottom panel, the orbital inclination angle $\iota$ continuously decreases over time, eventually bringing the SMO orbit nearly into coplanarity with the disk plane. In this work, the time for $\iota$ to reach the critical angle $\iota_{\rm crit}=3H_{\rm disk}/p_{\rm ini}$, is $t_{\rm cap}=5.81\times 10^{7}\ M_{\bullet}$.
% \begin{figure}
%     \centering
%     \includegraphics[scale=0.6]{figures/p_200_e_0_3_z1_40_deg/three_para_evo_p_200_e_0_3_z1_40_deg_t.pdf}
%     \caption{Evolution of the semi-major axis $a$, eccentricity $e$, and orbital inclination angle $\iota$ over coordinate time $t$, computed using the simplified model~\eqref{eq:simplified model}. Solid lines labeled ``orbit" show the results obtained with the method of extracting the parameters from the evolved E.O.M, while dashed lines labeled ``int ave" show the results obtained from adiabatic approximation. The zoomed-in panels display the variations of $a,\ e,\ \iota$ over the relatively short initial time interval $(0,\ 6\times 10^5\ M_{\bullet})$. In the bottom panel, the intersection point of the two red dotted lines indicates the critical inclination $\iota_{\rm crit}=3H_{\rm disk}/p_{\rm ini}$ and the corresponding time $t_{\rm cap}=5.81\times 10^7\ M_{\bullet}$.}
%     \label{fig:p 200 e 0.3 z1 40 deg three evo}
% \end{figure}

%\subsection{Parameter evolution with the aero-drag force model}



\subsection{Orbital evolution of stars}\label{sec:aerodynamic}
To assess the reliability of the adiabatic approximation in tracking secular evolution of a star colliding with an accretion disk, we first 
compare adiabatic evolution with full orbital evolution with the aero-drag force~\eqref{eq:aero drag}: the semi-major axis $a(t)\equiv p/(1-e^2)$, the eccentricity $e(t)$ and the inclination angle $\iota(t)$.  

% This does not contradict the results of Ref.~\cite{OConnor2020}, where a secular evolution of the apsidal angle was reported after averaging over an orbital period in the presence of thin-disk collisions. The difference originates from both the definition of the orbital phases and the use of the adiabatic approximation in this work. The phase offsets $\psi_{0}$ and $\chi_{0}$ are treated as integration constants specifying the initial orientation of the radial and polar motions, and their evolution is not tracked beyond leading adiabatic order.
%Consequently, information associated with per-orbit apsidal precession is not retained in the present framework. In contrast, Ref.~\cite{OConnor2020} explicitly constructs the apsidal angle by averaging the effects of thin-disk collisions over an orbital period, thereby capturing the cumulative pericenter advance from successive disk crossings. The secular evolution reported there thus reflects the long-term accumulation of per-orbit precession, while in the adiabatic treatment adopted here such precessional effects are averaged out and do not appear as a secular drift of the phase offsets.

As an example, we choose $\gamma_{0}=2\times 10^{-6}\ M_{\bullet}^{-1}$,  the initial conditions of the orbital parameters $p_{\rm ini}=200\ M_{\bullet},\ e_{\rm ini}=0.3,\ \iota_{\rm ini}=50^{\circ}$, and the initial phases  $\psi_{0\ \rm ini}=0$ and $\chi_{0\ \rm ini}=\pi/3$. The disk height to the midplane is set to $H_{\rm disk}=1.5\ M_{\bullet}$. As shown in Fig.~\ref{fig:aerodynamic p 200 e 0.3 z1 40 deg three evo}, the adiabatic approximation works well in tracking secular evolution of the stellar orbit.
The semi-major axis $a$ decays over time due to energy loss caused by star-disk collisions. Note that $a$ evolves slowly near the final stage when the star is close to being captured by the disk $\iota \rightarrow H_{\rm disk}/r$, because their relative speed and the drag force are greatly reduced. The eccentricity $e$ also decays with time and the orbit becomes
nearly circular with a final value $e_{\rm f}= 3.6\times 10^{-3}$ at the final stage of the evolution. Our numerical solution shows no secular evolution of the “initial” phases,  $\expval{\dot \psi_0} = \expval{\dot \chi_0} =0$, consistent with the expectation for the dissipative perturbing force as considered in this work~\cite{PhysRevD.83.044037}.

%In the bottom panel, the orbital inclination angle $\iota$ continuously decreases over time, eventually bringing the SMO orbit nearly into coplanarity with the disk plane. 
%In this work, the time for $\iota$ to reach the critical angle $\iota_{\rm crit}=3H_{\rm disk}/p_{\rm ini}$, is $t_{\rm cap}=2.73\times 10^{8}\ M_\bullet$. In this work, we also compute cases with $p_{\rm ini}=300\ M_{\bullet},\ 400\ M_{\bullet}$, while keeping all other initial parameters fixed. The corresponding capture times are $t_{\rm cap}=5.52\times 10^{8}\ M_{\bullet},\ 9.05\times 10^{8}\ M_{\bullet}$, respectively.\par
\begin{figure}
    \centering
    \includegraphics[scale=0.6]{figures/p_200_e_0_3_z1_40_v_rel/three_para_evo_p_200_e_0_3_z1_40_deg_v_rel_t.pdf}
    \caption{Evolution of the semi-major axis $a$, eccentricity $e$, and orbital inclination angle $\iota$ of a star. Solid lines labeled ``full" are obtained from full E.O.M. \eqref{eq: forced geodesic}, while dashed lines labeled ``adiabatic" show the secular evolution results obtained by evolving  Eq.~\eqref{eq:secular ode}. The zoomed-in panels display the variations of $a,\ e,\ \iota$ over the relatively short initial time interval $(0,\ 6\times 10^5\ M_{\bullet})$. In the bottom panel, the intersection point of the two red dotted lines indicates the critical inclination $\iota_{\rm crit}=3H_{\rm disk}/p_{\rm ini}$ and the corresponding time $t_{\rm cap}=2.7\times 10^{8}\ M_\bullet$.}
    \label{fig:aerodynamic p 200 e 0.3 z1 40 deg three evo}
\end{figure}


For the purpose of illustrating the eccentricity evolution driven by the aero-drag force, we compute the eccentricity evolution rate $\expval{\mathrm{d}e/\mathrm{d}\tau}$ as a function of the orbital inclination angle $\iota$, fixing the semi-latus rectum at $p = 300\ M_{\bullet}$ and considering representative eccentricity values $e = 0.1,\ 0.3,\ 0.7,\ 0.9$, as shown in Fig.~\ref{fig:edot value aero}. It is evident that star-disk collisions tend to circularize the star orbit for all inclination angles and orbital eccentricities.

%We also find that $\expval{\mathrm{d}e/\mathrm{d}\tau}$ exhibits a nonmonotonic dependence on the eccentricity $e$. This behavior arises from the relativistic osculating formulation, in which the secular evolution of $e$ results from the competition among radial, polar, and azimuthal drag components with distinct phase dependence and different scalings with $e$. As $e$ varies, the relative importance of these contributions changes, leading to a smooth but nonmonotonic response. {\bf Not much information in this explanation.}
%{\bf briefly explain why the dependence on $\iota$ is non-monotonic.}


\begin{figure}
    \centering
    \includegraphics[scale=0.55]{small_figures/ecc_dot_aero_p_300.pdf}
    \caption{The comparison of $\expval{\mathrm{d}e/\mathrm{d}\tau}$ variations with orbital inclination angle $\iota$ for different $e$, with $p=300\ M_{\bullet}$ assuming the aero-drag force model. The vertical black dotted line marks $\iota=\pi\ (180^{\circ})$.}
    \label{fig:edot value aero}
\end{figure}

To understand the final fate of a star colliding with a long-lived disk, we compare the disk-capture timescale $\tau_{\iota} := \abs{\iota/\expval{\dv{\iota}{\tau}}}$ and 
the orbital decay timescale $\tau_{a}:=\abs{a/\expval{\dv{a}{\tau}}}$ associated with the shrinkage of the semi–major axis $a$, where
\begin{equation}
    \expval{\dv{a}{\tau}}=\expval{\frac{1}{1-e^2}\dv{p}{\tau}+\frac{2ep}{(1-e^2)^2}\dv{e}{\tau}}.
\end{equation} 
As shown in Fig.~\ref{fig: timescale aero}, the capture timescale $\tau_{\iota}$ increases with increasing $\iota$, whereas the orbital decay timescale $\tau_{a}$ decreases due to the larger relative velocity between the star and the disk, which enhances the effect of the aero-drag force. As $\iota\rightarrow 0$, the timescale $\tau_{a}$ also decreases, however,  the star–disk collision model is no longer appropriate, since the orbit of the star becomes nearly embedded within the disk.

The comparison between $\tau_{\iota}$ and $\tau_{a}$ further shows that, for sufficiently large initial inclination angles, the orbital–decay timescale becomes shorter than the alignment timescale. This shows that a star on a high-inclination orbit will finally be disrupted by the  central SMBH before  being captured by the disk. 
On the other hand,  a star on a low-inclination orbit will be captured by the disk  before getting disrupted by the central SMBH.
However, we note that in the actual orbital evolution, the semi–major axis and eccentricity evolve simultaneously, which can modify the rates of change of the parameters. Consequently, such estimates based solely on the initial values may not fully coincide with the true evolutionary behavior. The above calculations also show that $\expval{\mathrm{d}{\iota}/\mathrm{d}{\tau}}<0$ and $\expval{\mathrm{d}{a}/\mathrm{d}{\tau}}<0$ over the entire range of $\iota$, consistent with the expectation and the actual evolution results found in this work.

% The timescale for the SMO orbit to align with the disk plane, approximated as the capture time, is characterized by the proper time $\tau_{\iota}$, estimated with $\abs{\iota_{\rm ini}/\expval{\dv{\iota}{\tau}}_{\rm ini}}$. Likewise, the orbital-decay timescale associated with the shrinkage of the semi–major axis $a$, is denoted by $\tau_{a}$, and estimated with $\abs{a_{\rm ini}/\expval{\dv{a}{\tau}}_{\rm ini}}$, where
% \begin{equation}
%     \expval{\dv{a}{\tau}}=\expval{\frac{1}{1-e^2}\dv{p}{\tau}+\frac{2ep}{(1-e^2)^2}\dv{e}{\tau}}.
% \end{equation} 
%These timescales are computed and presented 


\begin{figure}
    \centering
    \includegraphics[scale=0.5]{small_figures/aero_timescale_a_iota_contrast.pdf}
    \caption{The comparison of estimated shrink timescale with $\abs{a/\expval{\dv{a}{\tau}}}$ (solid dotted line) and capture timescale with $\abs{\iota/\expval{\dv{\iota}{\tau}}}$ (solid line) with orbital inclination angle $\iota$ for different $e$, fixing $p=300\ M_{\bullet}$, using the aero-drag force model for stars. The vertical black dotted line marks $\iota=\pi\ (180^{\circ})$.} %{\bf legends in two columns}}
    \label{fig: timescale aero}
\end{figure}
\par

%{\bf Logic connection between the last and the following paragraphs ?}
%After discussing the evolution behavior and timescales for the orbital parameters $a,\ e$ and $\iota$, here we will show several examples of the numerical calculation based on the adiabatic approximation. 
After establishing the evolution behavior and the associated timescales for the orbital parameters $a,\ e$ and $\iota$ under the aero-drag force, we now present several representative numerical examples of the orbital evolution computed within the adiabatic approximation.
We consider orbits with an initial semi-latus rectum $p_{\rm ini}=300\ M_{\bullet}$ and compute the evolution with different initial inclination angles, focusing on initial eccentricities of $e_{\rm ini}=0.3,\ 0.7$, under the aero-drag force model with $\gamma_{0}=2\times 10^{-6}\ M_{\bullet}^{-1}$.
The evolution of $a,\ e,\ \iota$ in the case of $e_{\rm ini}=0.3$ with different initial inclination angles under the aero-drag force model is shown in the three panels of Fig.~\ref{fig:aero e03 three evo}. 
%Results are presented for alignment $(\iota_{\rm ini}<90^{\circ})$ and anti-alignment cases $(\iota_{\rm ini}>90^{\circ})$.\par

As shown in the top and bottom panels of Fig.~\ref{fig:aero e03 three evo}, for the cases with $\iota_{\rm ini}\geq 130^{\circ}$, the orbits shrink, with the semi-major axis $a$ decreasing to below $10\ M_{\bullet}$, before alignment with the disk plane is achieved. For the cases with $\iota_{\rm ini}\leq 110^{\circ}$, the star successfully aligns with the disk plane and does not undergo significant shrinkage. The middle panel of Fig.~\ref{fig:aero e03 three evo} shows that the aero-drag force circularizes orbits with different inclination angles, regardless of the initial inclination.
We also find that stars on low-inclination orbits require longer capture times. This is primarily because, for low-inclination orbits, the relative velocity between the SMO and the disk material is small. According to Eq.~\eqref{eq:aero drag}, $\abs{\boldsymbol{F}_{\rm aero}}\propto \abs{\boldsymbol{v}_{\rm rel}}^{2}$, leading to a correspondingly slower orbital evolution. In addition, during the alignment process, the semi-major axis $a$ gradually decreases, which increases $\abs{\boldsymbol{v}_{\rm rel}}$ and accelerates the subsequent evolution. This effect further explains why low-inclination stars experience longer capture times: although high-inclination orbits require larger changes in inclination, they initially evolve more rapidly due to their larger relative velocities, and the accompanying shrinkage of $a$ during alignment further enhances this evolution, resulting in shorter capture times overall.
%We also find that stars on low-inclination orbits require longer capture times. It is because in the low-inclination orbits, the relative velocity between the SMO and disk material is small, and according to Eq.~\eqref{eq:aero drag}, $\abs{\boldsymbol{F}_{\rm aero}}\propto \abs{\boldsymbol{v}_{\rm rel}}^{2}$, the orbital evolution rate is thus smaller. Notice that the semi-major axis $a$ decreases during the alignment process, which will make $\abs{\boldsymbol{v}_{\rm rel}}$ larger, this also explains why this scenario happens.
%{\bf briefly explain why it takes longer time to capture a low-inclination star. }

The evolution of $a,\ e,\ \iota$ in the case of $e_{\rm ini}=0.7$ with the aero-drag force model is shown in the three panels of Fig.~\ref{fig:aero e07 three evo}, with all other conditions identical to those for $e_{\rm ini}=0.3$. Overall, changing the initial eccentricity does not significantly affect the evolution of the orbital parameters.
\begin{figure}
    \centering
    \includegraphics[scale=0.6]{Other_figures/three_para/three_para_evo_Aero_drag_model_contrast_e_0_3_r.pdf}
    \caption{Evolution of the semi-major axis $a$, eccentricity $e$, and orbital inclination angle $\iota$ using the aero-drag force model with a constant damping coefficient $\gamma_{0}=2\times 10^{-6}\ M_{\bullet}^{-1}$. The initial conditions are $p_{\rm ini}=300\ M_{\bullet}$, and $\ e_{\rm ini}=0.3$. The prograde cases $(\iota_{\rm ini}<90^{\circ})$ are shown with solid lines, while the retrograde cases $(\iota_{\rm ini}>90^{\rm \circ})$ are shown with dashed lines. The blue dotted horizontal line at $a=10\ M_{\bullet}$ in the top panel marks the threshold below which the orbit is considered to have shrunk in this work. The black dotted horizontal line in the bottom panel marks $\iota=\pi/2$.}
    %\caption{Semi-major axis $a$, eccentricity $e$, orbital inclination angle $\iota$ evolution with time, using the aero-drag force model with the constant part of damping coefficient $\gamma_{0}=2\times 10^{-6}\ M_{\bullet}^{-1}$. The initial conditions are $p_{\rm ini}=300\ M_{\bullet},\ e_{\rm ini}=0.3$. The alignment cases with $\iota_{\rm ini}<90^{\circ}$ are depicted with solid lines, and the anti-alignment cases with $\iota_{\rm ini}>90^{\rm \circ}$ are depicted with dashed lines. The blue dotted horizontal line is at $a=10\ M_{\bullet}$, indicating the threshold range within which the orbit is considered to have completed its decay in this work.}
    \label{fig:aero e03 three evo}
\end{figure}
\begin{figure}
    \centering
    \includegraphics[scale=0.6]{Other_figures/three_para/three_para_evo_Aero_drag_model_contrast_e_0_7_r.pdf}
    \caption{Same as Fig.~\ref{fig:aero e03 three evo} except for $e_{\rm ini}=0.7$.}
    \label{fig:aero e07 three evo}
\end{figure}
\subsection{Orbital evolution of sBHs}
%For the case of the SMO being an sBH, it is necessary to employ the dynamical friction model. In contrast to the aero-drag force model, it can lead to an increase in eccentricity within certain ranges of inclination angles, resulting in the orbital ellipticization. As shown in Fig.~\ref{fig:edot value sbh}, the time evolution rate of the eccentricity at different orbital inclination angles is calculated just like in Fig.~\ref{fig:edot value aero} except for using the dynamical model with $\gamma_{0}=4\times 10^{-13}\ M_{\bullet}^{-1}$. For a subset of initially aligned configurations, there is a tendency for the eccentricity to increase, whereas in all anti-aligned cases, the eccentricity inevitably grows. Only when the inclination is below a certain threshold does the dynamical friction model lead to orbital circularization behavior similar to that of the aero-drag model. For $e=0.1,\ 0.3,\ 0.7,\ 0.9$, the critical inclination angles for $\expval{\dv{e}{\tau}}$ to change value are $\iota_{\rm crit}=56^{\circ},\ 57.5^{\circ},\ 66.2^{\circ},\ 75.6^{\circ}$, respectively.\par
In this subsection, we consider  sBH–disk collisions, where the perturbing force is dominated by dynamical friction~\eqref{eq:Fdyn}. The change rate of the eccentricity $e$ as a function of the orbital inclination angle $\iota$ is computed using the same method and initial conditions as in Fig.~\ref{fig:edot value aero}, except that the dynamical-friction model with $\gamma_{0}=4\times10^{-13}\ M_{\bullet}^{-1}$ is adopted. The results are shown in Fig.~\ref{fig:edot value sbh}. 
In contrast to the aero-drag force model, we find dynamical friction can increase the orbital eccentricity $e$
if the orbiter is highly misaligned  with $\iota$ in the range of $(\iota_{\rm crit}, 180^\circ)$, where $\iota_{\rm crit} \approx 60^\circ - 90^\circ$ for the parameter range considered in this work. Similar results were also found in previous studies \cite{Wang2024, spieksma2025gripdiskdraggingcompanion}, which motivate some recent discussions on eccentric wet EMRIs. As we will show later in this section, EMRIs that are captured by disks are circularized due to the eccentricity damping when $\iota < \iota_{\rm crit}$.

%Only for inclinations below a certain threshold does the dynamical friction model drive orbital circularization similar to that led by the aero–drag force model. 
%For $e=0.1,\ 0.3,\ 0.7,\ 0.9$, the critical inclination angles at which $\expval{\mathrm{d}e/\mathrm{d}\tau}$ changes sign are $\iota_{\rm crit}=56^{\circ},\ 57.5^{\circ},\ 66.2^{\circ},\ 75.6^{\circ}$, respectively. 
%The eccentricity excitation predicted by the dynamical friction model in this work is consistent with the findings of Refs.~\cite{Wang2024, spieksma2025gripdiskdraggingcompanion}, as well as with the orbital evolution behavior obtained in our numerical examples presented in Figs.~\ref{fig:friction e03 three evo} and~\ref{fig:friction e07 three evo}.
%The eccentricity excitation predicted by the dynamical friction model in this work is consistent with the results reported in Refs.~\cite{Wang2024, spieksma2025gripdiskdraggingcompanion}, and also the evolution results obtained in this work, which are shown in Section~\ref{sec:evolution examples}.
%{\bf Explain why $e$ grows }

%Just like in Sec.~\ref{sec:aerodynamic}, two timescales for the sBH orbit to align with the disk plane or shrink are calculated and shown in Fig.~\ref{fig: timescale sbh}, with $p=300\ M_{\bullet}$, and the same set of eccentricity values as in Fig.~\ref{fig:edot value sbh}, there are also $\expval{\dv{a}{\tau}}<0$ and $\expval{\dv{\iota}{\tau}}<0$ according to the calculation results.\par
As discussed in Sec.~\ref{sec:aerodynamic}, we also calculate two timescales $\tau_{a}$ and $\tau_{\iota}$ for the sBH orbit, as shown in Fig.~\ref{fig: timescale sbh}, with the same set of initial conditions as in Fig.~\ref{fig:edot value sbh}. The calculations also yield $\expval{\mathrm{d}a/\mathrm{d}\tau}<0$ and $\expval{\mathrm{d}\iota/\mathrm{d}\tau}<0$ for all inclination angles, showing that the disk is aligning the orbiter. 
From Fig.~\ref{fig: timescale sbh}, for the sBH-SMO cases, large initial orbital inclination angles result in $\tau_{a}<\tau_{\iota}$, consistent with the behavior found in the stellar-SMO cases, while $\tau_{\iota}$ increases monotonically with $\iota$, $\tau_{a}$ initially increases with $\iota$ because the larger relative velocity between the sBH and the disk material reduces the effect of the dynamical friction according to Eqs.~\eqref{eq:Fdyn} and~\eqref{eq:dynamical friction gamma}, in contrast to the behavior under the aero–drag force model. It can also be found that $\tau_{a}$ decreases as $\iota\rightarrow\pi$, which can be interpreted as the strengthened interaction resulting from the sBH orbit gradually becoming embedded in the disk plane.
\begin{figure}
    \centering
    \includegraphics[scale=0.55]{small_figures/sBH_p_300_ecc_dot.pdf}
    \caption{The comparison of $\expval{\dv{e}{\tau}}$ variations with orbital inclination angle $\iota$ for different $e$, with $p=300\ M_{\bullet}$, with the dynamical friction model for the sBH-SMO cases. The vertical black dotted line marks $\iota=\pi\ (180^{\circ})$.}
    \label{fig:edot value sbh}
\end{figure}
\begin{figure}
    \centering
    \includegraphics[scale=0.5]{small_figures/sBH_timescale_a_iota_contrast.pdf}
    \caption{The comparison of estimated shrink timescale with $\tau_{a}=\abs{a/\expval{\dv{a}{\tau}}}$ (solid dotted line) and capture timescale with $\tau_{\iota}=\abs{\iota/\expval{\dv{\iota}{\tau}}}$ (solid line) with orbital inclination angle $\iota$ for different $e$, fixing $p=300\ M_{\bullet}$, using the dynamical friction force model for the sBH-SMO cases. The vertical black dotted line marks $\iota=\pi\ (180^{\circ})$.}
    \label{fig: timescale sbh}
\end{figure}
\par
To illustrate the evolution behavior and timescale hierarchy discussed above, we now present representative numerical examples of the orbital evolution under the dynamical friction, computed within the adiabatic approximation. We consider orbits with an initial semi-latus rectum $p_{\rm ini}=300\ M_{\bullet}$ and compute the evolution with different initial inclination angles, focusing on initial eccentricities of $e_{\rm ini}=0.3,\ 0.7$, under the dynamical-friction model with $\gamma_{0}=4\times 10^{-13}\ M_{\bullet}^{-1}$.
The evolution of $a,\ e,\ \iota$ for $e_{\rm ini}=0.3$ under the dynamical friction model is shown in the three panels of Fig.~\ref{fig:friction e03 three evo}. In this work, for dynamical friction model, the calculation is terminated once the orbit approaches the near-aligned configuration with the disk plane, as during the alignment process, the relative velocity between the sBH and the disk  gradually decreases, the change rate of the orbital parameters becomes sufficiently large and potentially invalidates the adiabatic approximation.

For cases with $\iota_{\rm ini}=20^{\circ}-135^{\circ}$, near-complete alignment of the orbits occurs before significant orbital shrinkage. In contrast, for the higher inclination case with $\iota_{\rm ini}=170^{\circ}$, although the semi-major axis $a$ is still several times of $10\ M_{\bullet}$, the calculation is terminated before the condition $p-6-2e>0$ is violated~\cite{PhysRevD.50.3816}, which marks the transition from a stable to an unstable orbit and the inevitable infall of the sBH into the central SMBH.\par
As shown in the middle panel of Fig.~\ref{fig:friction e03 three evo}, for all retrograde cases, the eccentricity is excited (see Fig.~\ref{fig:edot value sbh}). For prograde cases, exemplified by $\iota_{\rm ini}=90^{\circ}$, the eccentricity is also initially excited. As $\iota$ decreases, a transition occurs from $\expval{\mathrm{d}e/\mathrm{d}\tau}>0$ to $\expval{\mathrm{d}e/\mathrm{d}\tau}<0$, and the eccentricity eventually declines. For smaller initial inclination angles, such as $\iota_{\rm ini}=20^{\circ},\ 45^{\circ}$, the eccentricity decreases from the outset.\par
The evolution of $a,\ e,\ \iota$ for $e_{\rm ini}=0.7$ with the dynamical friction model is shown in the three panels of Fig.~\ref{fig:friction e07 three evo}, with all other conditions identical to those for $e_{\rm ini}=0.3$. Overall, the evolution of the orbital parameters shows no significant differences compared to the $e_{\rm ini}=0.3$ case.%, except for a slightly shorter evolution timescale.
\begin{figure}
    \centering
    \includegraphics[scale=0.6]{Other_figures/three_para/three_para_evo_sBH_contrast_e_0_3_r.pdf}
    \caption{Same as Fig.~\ref{fig:aero e03 three evo} except for using dynamical friction model with the constant part of damping coefficient $\gamma_{0}=4\times 10^{-13}\ M_{\bullet}^{-1}$. }
    \label{fig:friction e03 three evo}
\end{figure}
\begin{figure}
    \centering
    \includegraphics[scale=0.6]{Other_figures/three_para/three_para_evo_sBH_contrast_e_0_7_r.pdf}
    \caption{Same as Fig.~\ref{fig:friction e03 three evo} except for $e_{\rm ini}=0.7$.}
    \label{fig:friction e07 three evo}
\end{figure}


\subsection{Capture timescales of AGN disks}\label{sec:capture timescales}
%{\bf merge your [Dependence of the capture timescale on $p$] and [Estimation for the capture timescale] in this subsection}
As shown in the preceding subsections, collisions with the disk drive a secular decrease in the orbital inclination angle $\iota$, which may ultimately lead to the SMO being captured by the disk. In this subsection, we estimate characteristic timescales for SMOs to be captured by the disk in EMRI systems over a wide range of orbital scales. 

We first analyze the dependence of the capture timescale $t_{\rm cap}$ on relevant quantities of the EMRI+disk system.
For clarity, starting from the aero-drag force model in Eq.~\eqref{eq:aero-drag force gamma}, we note that the coefficients $C_{z_1}(\theta),\ C_{z_{1}}(\phi)$ given in Appendix~\ref{appendix:coefficients} scale as $p^{3/2}$, while the $\theta,\ \phi$ acceleration components in Appendix~\ref{appendix:form of acceleration} scale as
\begin{equation}
    \gamma_{0}\ \abs{\boldsymbol{v}_{\rm rel}}\ p^{-3/2}\sim \gamma_{0}\ p^{-1/2}\ p^{-3/2}\sim \gamma_{0}\ p^{-2},
\end{equation}
 where $\abs{\boldsymbol{v}_{\rm rel}}$ can be approximated using the Keplerian relation that $\abs{\boldsymbol{v}_{\rm rel}}\sim p^{-1/2}$. Combining this with the crossing time of the SMO through the disk with the height $H_{\rm disk}$, which can be approximated using the Keplerian estimate $t_{\rm cross}\sim H_{\rm disk}\ v_{\rm SMO}^{-1}\sim H_{\rm disk}\ p^{1/2}$, and using $z_{1}=\sin\iota$, we find that the change in $\iota$ produced by each SMO–disk collision scales as
\begin{equation}
    \Delta\iota\sim\Delta z_{1}\sim p^{\frac{3}{2}}\cdot \gamma_{0}\ p^{-2}\cdot H_{\rm disk}\ p^{\frac{1}{2}}\sim \Sigma_{\rm g}R_{\star}^2m_{\star}^{-1}.
\end{equation}
Thus, the total number of collisions required for a given inclination angle $\iota$ to decrease from its initial value to nearly zero scales as
\begin{equation}
    N_{\rm cap}\sim \Sigma_{\rm g}^{-1}R_{\star}^{-2}m_{\star}.
\end{equation}
Since the orbital period of the SMO scales as $T_{\rm obt}\sim p^{3/2}$,  the capture timescale follows
\begin{equation}
    t_{\rm cap}\sim N_{\rm cap}\ T_{\rm obt}\sim \Sigma_{\rm g}^{-1}R_{\star}^{-2}m_{\star} p^{3/2}.
\end{equation}
The power law relations $t_{\rm cap}\sim p^{n}$ for the dynamical friction model can be derived in a similar manner. In both cases, the dependence of the damping coefficient on the relative velocity $\abs{\boldsymbol{v}_{\rm rel}}$ should be taken into account. The resulting power law relations for both interaction models are summarized in Table~\ref{tab:power tau p}.\par
As a numerical example, shown in Fig.~\ref{fig:time_scale p list}, we examine the dependence of the capture timescale $t_{\rm cap}$ on the initial semi-latus rectum $p_{\rm ini}$ using the aero-drag force model with the damping coefficient $\gamma_{0}=2\times 10^{-6}\ M_{\bullet}^{-1}$. $p_{\rm ini}$ is sampled over in the range $(300,\ 2000)\ M_{\bullet}$, for two initial inclination angles, $\iota_{\rm ini}=45^{\circ}$ and $90^{\circ}$, while the initial eccentricity fixed at $e_{\rm ini}=0.3$. For clarity, we define the capture timescale $t_{\rm cap}$ as the coordinate time at which the orbital inclination decreases from its initial value $\iota_{\rm ini}$ to a small threshold $\iota_{\rm crit}=1.5\times10^{-2}$.
The numerical results are well consistent with the expected scaling  $t_{\rm cap}\sim p^{3/2}$. 


\begin{table}[htbp]
  \caption{Power scaling relations between $t_{\rm cap}$ and $p$ and other parameters}
  \label{tab:power tau p}
  \begin{tabular}{l|ccc}
    %\hline
    Model & $f^{\mu}$ & $N_{\rm cap}$ & $t_{\rm cap}$ \\
    \hline
    Aero-drag & $\sim \gamma_{0}\abs{\bm{v}_{\rm vel}}u^{\mu}_{\rm vel}$ & $\sim \Sigma_{\rm g}^{-1} R_{\star}^{-2}m_{\star}p^{0}$ & $\sim \Sigma_{\rm g}^{-1} R_{\star}^{-2}m_{\star}p^{\frac{3}{2}}$ \\
    Dynamical friction & $\sim \gamma_{0}\abs{\bm{v}_{\rm vel}}^{-3}u^{\mu}_{\rm vel}$ & $\sim\Sigma_{\rm g}^{-1} m_{\bullet}^{-1} p^{-2}$ & $\sim \Sigma_{\rm g}^{-1} m_{\bullet}^{-1}p^{-\frac{1}{2}}$ \\
    %\hline
  \end{tabular}
\end{table}
\begin{figure}
    \centering
    \includegraphics[scale=0.55]{figures/t_cap_p_aero_drag.pdf}
    \caption{Dependence of the capture timescale on the initial semi-latus rectum $p_{\rm ini}$ over the range $300\ M_{\bullet}$ to $2000\ M_{\bullet}$, computed using the aero-drag force model. }
    \label{fig:time_scale p list}
\end{figure}
\label{sec:discussion}
\par
To assess whether misaligned EMRIs can be captured by the AGN disk within a certain timescale, we estimate the disk-capture timescale using scaling relations in Table~\ref{tab:power tau p}. The capture timescale is evaluated with $p_{\rm ini}=300\ M_{\bullet}$, and $e_{\rm ini}=0.3$, in combination with the Sirko \& Goodman (SG) disk model~\cite{10.1046/j.1365-8711.2003.06431.x}.
For simplicity, we restrict our analysis to cases in which orbital alignment occurs prior to significant orbital shrinkage. 
We adopt the same conventions as in Ref.~\cite{Wang2024} for the SMO mass, taking $m_{\star}=m_{\bullet}=30\ M_{\odot}$, and for stellar SMOs, the stellar radius is set to $R_{\star}=R_{\odot}(m_{\star}/M_{\odot})^{\frac{3}{4}}$. %{[\bf Use $R_\star, m_\star$]}
The disk surface-density and scale-height profiles are obtained using the \texttt{Sirko} module in the \texttt{pagn} package ~\cite{Gangardt_2024} (see Fig.~\ref{fig:Sigmag Hdisk pagn}).

\begin{figure}[htbp]
    \centering
    \includegraphics[scale=0.55]{pagn_models/Sigma_g_H_disk_SMBH_mass_contrast.pdf}
    \caption{The surface density $\Sigma_{\rm g}$ and disk height $H_{\rm disk}$ distribution with the radial distance $r$ from the central SMBH with different masses, for SG disk model, with the Eddington ratio $l_{e}=0.5$, viscosity parameter $\alpha=0.1$, hydrogen abundance $X=0.7$, and $\alpha-$disk case with $b=0$. }
    \label{fig:Sigmag Hdisk pagn}
\end{figure}


For the aero-drag force model, the results presented in Figs.~\ref{fig:aero e03 three evo} and~\ref{fig:aero e07 three evo} show that the capture timescale in prograde cases is longer than  in retrograde cases.
The estimates are carried out for initial inclination angles $\iota_{\rm ini}=10^{\circ},\ 30^{\circ},\ 50^{\circ},\ 70^{\circ},\ 90^{\circ}$ and for SMBH masses of $M_{\bullet}=10^6\ M_{\odot}, 10^{8}\ M_{\odot}$, with the results shown in two panels of Fig.~\ref{fig:aero estimation 1e6 1e8}, respectively.
For prograde cases under the aero-drag force model, the capture timescale exhibits only a weak dependence on the initial orbital inclination angle, as can be seen in Figs.~\ref{fig:aero e03 three evo} and \ref{fig:aero e07 three evo}. This behavior can be interpreted as follows. Larger orbital inclinations lead to higher relative velocities, which in turn produce larger inclination-change rates. As a result, once the inclination exceeds a certain threshold, the alignment timescale decreases.\par
\begin{figure}[htbp]
    \centering
    \includegraphics[scale=0.55]{small_figures/aero_cap_time_estimation_1e6_alpha_0_1.pdf}\\
    \includegraphics[scale=0.55]{small_figures/aero_cap_time_estimation_1e8_alpha_0_1.pdf} 
    \caption{Capture timescales estimated using the aero-drag force model for initial inclination angles $\iota_{\rm ini}=10^{\circ},\ 30^{\circ},\ 50^{\circ},\ 70^{\circ},\ 90^{\circ}$. The horizontal red dashed line indicates $t=10^{6}\ \mathrm{yr}$. The top and bottom panels correspond to SMBH masses of $M_{\bullet}=10^{6}\ M_{\odot}$ and $10^{8}\ M_{\odot}$, respectively.} %{\bf Complete the fig description}
    \label{fig:aero estimation 1e6 1e8}
\end{figure}
% \begin{figure}[htbp]
%     \centering
%     \includegraphics[scale=0.55]{small_figures/aero_cap_time_estimation_1e8.pdf}
%     \caption{The estimation for the alignment timescale with $M_{\bullet}=10^8\ M_{\odot}$ with the aero-drag force model. The horizontal red dashed line corresponds to $t=10^6$ years.}
% \end{figure}
%For the dynamical-friction force model, the timescale estimation is studied with the SMBH masses of $M_{\bullet}=10^6\ M_{\odot},\ 10^{7}\ M_{\odot}$ and $10^8\ M_{\odot}$, as shown in Fig.~\ref{fig:sbh estimation 1e6}, Fig.~\ref{fig:sbh estimation 1e7} and Fig.~\ref{fig:sbh estimation 1e8}, respectively.
For the dynamical friction model, the timescale estimates are evaluated for SMBH masses of $M_{\bullet}=10^6\ M_{\odot},\ 10^{7}\ M_{\odot}$ and $10^8\ M_{\odot}$, as shown in three panels of Fig.~\ref{fig:sbh estimation 1e6 1e7 1e8}, respectively.
In the dynamical friction model, the capture timescale increases significantly with orbital inclination, since larger relative velocities lead to smaller inclination-change rates.\par
\begin{figure}[htbp]
    \centering
    \includegraphics[scale=0.55]{small_figures/sBH_cap_time_estimation_1e6_alpha_0_1.pdf}\\
     \includegraphics[scale=0.55]{small_figures/sBH_cap_time_estimation_1e7_alpha_0_1.pdf}\\
      \includegraphics[scale=0.55]{small_figures/sBH_cap_time_estimation_1e8_alpha_0_1.pdf}
    \caption{Capture timescales estimated using the dynamical friction model for initial inclination angles $\iota_{\rm ini}=10^{\circ},\ 20^{\circ},\ 30^{\circ},\ 45^{\circ},\ 60^{\circ},\ 90^{\circ},\ 105^{\circ},\ 120^{\circ},\ 135^{\circ}$. The horizontal red dashed line indicates $t=10^{6}\ \mathrm{yr}$. The top, middle and bottom panels correspond to SMBH masses of $M_{\bullet}=10^{6}\ M_{\odot},\ 10^{7}\ M_{\odot}$ and $10^{8}\ M_{\odot}$, respectively.}
    \label{fig:sbh estimation 1e6 1e7 1e8}
\end{figure}
% \begin{figure}[htbp]
%     \centering
%     \includegraphics[scale=0.55]{small_figures/sBH_cap_time_estimation_1e7.pdf}
%     \caption{The estimation for the alignment timescale with $M_{\bullet}=10^7\ M_{\odot}$ with the dynamical-friction force model. The horizontal red dashed line corresponds to $t=10^6$ years.}
%     \label{fig:sbh estimation 1e7}
% \end{figure}
% \begin{figure}[htbp]
%     \centering
%     \includegraphics[scale=0.55]{small_figures/sBH_cap_time_estimation_1e8.pdf}
%     \caption{The estimation for the alignment timescale with $M_{\bullet}=10^8\ M_{\odot}$ with the dynamical-friction force model. The horizontal red dashed line corresponds to $t=10^6$ years.}
%     \label{fig:sbh estimation 1e8}
% \end{figure}
For both models, the capture timescale is compared with a critical timescale of $10^{6}$ yr (1 Myr).
For the aero-drag force model, SMOs at distances of order $10^{-3}\ \rm pc$ can be captured within 1 Myr for $M_{\bullet}=10^6\ M_{\odot}$, whereas this distance increases to about $10^{-2}\ \rm pc$ for $M_{\bullet}=10^{8}\ M_{\odot}$.
For the dynamical friction model, the range of orbital inclination angles for which SMOs can be captured within 1 Myr decreases noticeably with increasing SMBH mass, primarily due to the disk surface-density profile and the scaling between the SMBH mass $M_{\bullet}$ and physical time. 
%The distribution profiles of the estimated capture timescales for $M_{\bullet}=10^{6}\ M_{\odot}$ differ from those in the other two cases, primarily due to differences in the disk surface density profiles (see Fig.~\ref{fig:Sigmag Hdisk pagn}).
%We find that the distribution profiles of the estimated capture timescales for the case of $M_{\bullet}=10^{6}\ M_{\odot}$ are different from the other two cases, which is mainly because of the discrepancy of the disk surface density profiles (see Fig.~\ref{fig:Sigmag Hdisk pagn}).
%\bf briefly explain why 1e6 looks different}

\subsection{Comparison with previous studies}\label{sec:comparison}
%{\bf Also compare https://iopscience.iop.org/article/10.3847/2041-8213/abe11d/pdf }
In previous studies~\cite{spieksma2025gripdiskdraggingcompanion,Wang2024,Rowan_2025,10.1093/mnras/staa3004,10.1093/mnras/stad1295,Secunda_2021,MacLeod:2019jxd,OConnor2020} of star–disk collisions in EMRIs, the orbital evolution of the SMO is typically modeled within a Newtonian framework. The secular changes are commonly estimated by summing the contributions from the two disk crossings within one orbital period and dividing by the orbital period. In this section, we compare the evolution of the orbital parameters $a,\ e$ and $\iota$ for the stellar-SMO and sBH-SMO cases, obtained under the adiabatic approximation in Schwarzschild spacetime and shown in Figs.~\ref{fig:aero e03 three evo}–\ref{fig:friction e07 three evo}, with those reported in previous studies.\par

Our calculations broadly agree with the results of \cite{Wang2024}.
Consistent with \cite{spieksma2025gripdiskdraggingcompanion} and \cite{Wang2024}, we find that collisions with the disk always tend to align the SMO no matter what the initial orbital inclination $\iota_{\rm ini}$ is.
An opposite trend was claimed in \cite{10.1093/mnras/stad1295} for sBHs with large initial inclination angles. 
%Within their framework, this behavior may be attributed to the direction of the drag-induced torque exerted by the disk on the sBH: when the inclination is sufficiently large, the geometric configuration of repeated disk crossings causes the torque to act in a way that drives the orbit toward larger inclination angles, as inferred from their geometric model. 
The authors cautioned that this peculiar behavior might reflect missing physics, such as the inconsistent assumption of circular orbits and the neglect of disk back reaction on the orbiting object.
Consistent with \cite{spieksma2025gripdiskdraggingcompanion} and \cite{Wang2024} , we find star-disk collisions always circularize the star orbit, while sBH-disk collisions excite the orbital eccentricity when the orbital inclination $\iota$ is large. But the temporary orbital eccentricity excitation does not lead to formation of eccentric wet EMRIs because the orbital eccentricity is largely damped when the inclination decreases and approaches zero when the orbiter is captured by the disk (Figs.~\ref{fig:friction e03 three evo}, \ref{fig:friction e07 three evo}). 

In \cite{Secunda_2021},  the authors found that the orbital eccentricity of embedded retrograde EMRIs can be efficiently excited to near unity by dynamical friction. This behavior is qualitatively consistent with the eccentricity growth observed in our high-inclination sBH cases, particularly for $\iota_{\rm ini}=170^{\circ}$ (see Figs.~\ref{fig:friction e03 three evo}, \ref{fig:friction e07 three evo}).
The finding has been used for argument of formation of eccentric wet EMRIs in AGN disks. In fact, as we have found in this work by tracking full orbital evolution of EMRIs,  AGN disks always tend to align EMRIs that are initially misaligned, and the captured EMRIs are effectively circularized though the orbital eccentricity may be temporarily excited during the capture process (Figs.~\ref{fig:friction e03 three evo}, \ref{fig:friction e07 three evo}).  As a result, there is little parameter space for forming eccentric wet retrograde EMRIs. 
\par
In addition to qualitative agreement, our calculations also show quantitative agreement with \cite{Wang2024} e.g., in the capture timescale of sBHs by AGN disks. 
For the sBHs, we adopt a constant damping coefficient $\gamma_{0}=4\times 10^{-13}\ M_{\bullet}^{-1}$, together with the disk scale height $H_{\rm disk}=1.5\ M_{\bullet}$. According to Eq.~\eqref{eq:dynamical friction gamma}, the remaining factors entering the damping coefficient can be combined into the dimensionless normalization $(\Sigma_{\rm g}/10^5\ \mathrm{g\cdot cm^{-2}})\cdot(m_{\bullet}/M_{\odot})\approx 86.58$. In Ref.~\cite{Wang2024}, sBHs were modeled with an initial semi-major axis $a_{\rm ini}=10^{-3}\ \mathrm{pc}$ and eccentricity $e_{\rm ini}=0.7$, which corresponds to an initial semi-latus rectum $p_{\rm ini}\sim 10^2\ M_{\bullet}$, for a central SMBH of $M_{\bullet}=10^{8}\ M_{\odot}$.
Adopting the surface density profile implied by the $\alpha-$disk model shown in their work, and using the scaling relations listed in Table~\ref{tab:power tau p}, together with the numerical results obtained in this work (see Fig.~\ref{fig:friction e07 three evo}), we find that the resulting capture timescales are consistent at the order-of-magnitude level. We note that the capture timescales reported in Ref.~\cite{Wang2024} are systematically shorter than those obtained here, likely because their disk model implies an increasing surface density as the orbit shrinks, whereas a constant surface density assumption is adopted in this work; nevertheless, the overall scaling remains consistent.
%For the sBH-SMO cases, we have adopted $\gamma_{0}=4\times 10^{-13}\ M_{\bullet}^{-1}$, with the same disk scale height $H_{\rm disk}=1.5\ M_{\bullet}$. According to Eq.~\eqref{eq:dynamical friction gamma}, the remaining factors contributing to the damping coefficient are combined into $\Sigma_{5}(m_{\bullet}/M_{\odot})\approx 86.58$.
%In \cite{Wang2024}, the sBH-SMO systems were modeled with the initial semi-major axis $a_{\rm ini}=10^{-3}\ \mathrm{pc}$ and eccentricity $e_{\rm ini}=0.7$, corresponding to $p_{\rm ini}\sim 10^2\ M_{\bullet}$ for $M_{\bullet}=10^{8}\ M_{\odot}$. If we take the surface density profile given by $\alpha-$ disk model shown in their work, according to the scaling relation in Table.~\ref{tab:power tau p} and the results obtained in this work (see Fig.~\ref{fig:friction e07 three evo}), The estimated capture timescales are of the same order of magnitude, although the capture timescales in \cite{Wang2024} are shorter, which may be the reason that as the orbits shrink, there will be larger surface density but not a constant assumed in this work.
%{[\bf one example is sufficient]}



%\section{Discussion and Conclusion}
%{\bf Move your [Comparison with other studies] here}


% ~\\


% xxxx

% z
% Different choices of the SMO mass and size, as well as disk surface density models, have been adopted in other previous studies, resulting in different effective damping coefficients.
% In this work, for the stellar-SMO cases, the damping coefficient is set to $\gamma_{0}=2\times 10^{-6}\ M_{\bullet}^{-1}$, with a disk scale height $H_{\rm disk}=1.5\ M_{\bullet}$. According to Eq.~\eqref{eq:aero-drag force gamma}, the remaining factors contributing to the damping coefficient are combined into $\Sigma_{5}(R_{\star}/R_{\odot})^2(m_{\star}/M_{\odot})^{-1}\approx 78.43$. 
% For the sBH-SMO cases, we have adopted $\gamma_{0}=4\times 10^{-13}\ M_{\bullet}^{-1}$, with the same disk scale height $H_{\rm disk}=1.5\ M_{\bullet}$. According to Eq.~\eqref{eq:dynamical friction gamma}, the remaining factors contributing to the damping coefficient are combined into $\Sigma_{5}(m_{\bullet}/M_{\odot})\approx 86.58$. These combined parameter values entering the damping coefficient are used to estimate the capture timescales following the method adopted in this work, and to facilitate a comparison with results from previous studies.\par
% We first focus on comparing the evolutionary behavior of the orbital parameters. In this work, we find that both the semi-major axis $a$ and the inclination angle $\iota$ decrease monotonically during the star–disk interaction for both the aero-drag and dynamical friction models. This behavior is consistent with the results obtained using the dynamical friction model reported in Ref.~\cite{spieksma2025gripdiskdraggingcompanion}, as well as with the $N$-body simulations incorporating both models presented in Ref.~\cite{Wang2024}. In addition, the eccentricity excitation observed in the dynamical friction model in this work has also been reported in Refs.~\cite{spieksma2025gripdiskdraggingcompanion,Wang2024}. However, in Ref.~\cite{10.1093/mnras/stad1295}, it was reported that, for sBH–SMO systems with the initial inclination angles $\iota_{\rm ini}>90^{\circ}$, $\iota$ increases once it exceeds a critical value. This behavior is opposite to the inclination-damping trend found in this work and in most previous studies. Such discrepancy was attributed to the assumption of circular orbits or to the neglect of additional physical effects {\bf be more clear: circular orbits lead to anti-alignment ??? what is addiation phyiscal effect?}.\par

% Secondly, we focus on the comparison of the capture timescale, an relate our results to several previous studies.
% \begin{enumerate}
%     \item In the study on sBH-SMO systems in Ref.~\cite{spieksma2025gripdiskdraggingcompanion}, the initial semi-latus rectum was taken to be $p_{\rm ini}\sim 10^5\ M_{\bullet}$, with an sBH mass $m_{\bullet}=10^{3}\ M_{\odot}$ orbiting an SMBH of mass $\ M_{\bullet}=10^7\ M_{\odot}$, Adopting the Sirko–Goodman (SG) disk model~\cite{10.1046/j.1365-8711.2003.06431.x}, these parameters correspond to $\Sigma_{5}(m_{\bullet}/M_{\odot})\approx 10$. Using the scaling relations summarized in Table~\ref{tab:power tau p}, and comparing with the value $\Sigma_{5}(m_{\bullet}/M_{\odot})\approx 86.58$ in this work, we find that the capture orbital period number $N_{\rm cap}\sim 10^3$ reported in Ref.~\cite{spieksma2025gripdiskdraggingcompanion} is in reasonable agreement with the results obtained in this work.
%     \item According to the $N$-body simulation results with an SMO the mass of $30\ M_{\odot}$ orbiting an SMBH of mass $M_{\bullet}=10^{8}\ M_{\odot}$ in Ref.~\cite{Wang2024}, the stellar–SMO system was modeled with the initial semi-latus rectum taken to be $p_{\rm ini}\sim 10^3\ M_{\bullet}$, and the stellar radius set to $R_{\star}=(m_{\star}/M_{\odot})^{\frac{3}{4}}R_{\odot}$. If the $\alpha-$disk model with $\alpha=0.01$ was adopted, using the scaling relation and also the combined value $\Sigma_{5}(R_{\star}/R_{\odot})^2(m_{\star}/M_{\odot})^{-1}\approx 78.43$ obtained in this work, we find that the estimated capture timescale may be $5$–$10$ times longer than those reported in Ref.~\cite{Wang2024}. For the sBH-SMO system modeled with $p_{\rm ini}\sim 10^2\ M_{\bullet}$, we use the scaling relation combined with the calculation results of the dynamical friction model presented in this work. The estimated capture timescales are of the same order of magnitude as those reported in Ref.~\cite{Wang2024}. 
%     \item In Ref.~\cite{Rowan_2025}, the evolution of an sBH–SMO system was studied for a stellar-mass black hole of mass $m_{\bullet}=25\ M_{\odot}$ circularly orbiting an SMBH of mass $M_{\bullet}=10^7\ M_{\odot}$, with the SMO initially located at an orbital radius of $\sim 10^4\ M_{\bullet}$. For the estimated disk surface density of $\Sigma_{\rm g}\approx 1.6\times 10^4\ \mathrm{g\ cm^{-2}}$ given in Ref.~\cite{Rowan_2025}, the alignment timescale for the largest initial inclination considered, $\iota_{\rm ini}=15^{\circ}$, was $10^5$ yr. This timescale is consistent with our calculation-based estimate and the scaling relations derived in this work.
%     \item In Ref.~\cite{10.1093/mnras/staa3004}, the capture of SMOs was studied under the assumption of circular orbits. Prograde configurations were considered for several types of stellar-SMOs orbiting a central SMBH of mass $M_{\bullet}=10^{8}\ M_{\odot}$. Here, We compare the results obtained from the SG disk model with those presented in this work. For the  stellar-SMO systems studied in Ref.~\cite{10.1093/mnras/staa3004}, where the SMOs were initially located at radii of $10^2-10^3\ R_{\rm g}$, the capture timescales reported for different stellar types are consistent, at the order-of-magnitude level, with the estimates derived from the scaling relations listed in this work.
% \end{enumerate}
% There are also studies of EMRI+disk systems in which the SMO orbits within the disk plane, such as Ref.~\cite{Secunda_2021}. In that work, the orbital evolution of an sBH embedded in an AGN disk on a retrograde orbit was investigated. They found that the eccentricity of the sBH can be excited to values close to unity, accompanied by a damping of the semi-major axis. This behavior is qualitatively similar to the evolution observed in our low-inclination retrograde sBH-SMO cases with an initial inclination of $\iota_{\rm ini}=170^{\circ}$. The evolution timescales reported in Ref.~\cite{Secunda_2021} are, however, significantly shorter than those found in this work. This difference arises because, in their setup, the sBH continuously interacts with the disk material along the entire orbit, whereas in our model the sBH interacts with the disk only through discrete disk-crossing events.
%In Ref.~\cite{spieksma2025gripdiskdraggingcompanion}, the orbital evolution of an EMRI system composed of a central SMBH with mass $M_{\bullet}=10^{7}\ M_{\odot}$ and an sBH-SMO with mass $m_{\bullet}=10^{3}\ M_{\odot}$ was investigated within a Newtonian framework, adopting the dynamical friction model. The initial semi-major axis was set to $a_{\rm ini}=10^{6}\ M_{\bullet}$, corresponding to an initial semi-latus rectum $p_{\rm ini}\sim 10^5\ M_{\bullet}$ through the relation $p=a(1-e^2)$. Among the disk models considered, the Sirko–Goodman (SG) disk~\cite{10.1046/j.1365-8711.2003.06431.x} was adopted as a representative case. As shown in Fig.~\ref{fig:Sigmag pagn}, for $M_{\bullet}=10^7\ M_{\odot}$, the surface density on spatial scales of $\sim 10^{5}$–$10^{6}\ M_{\bullet}$ is $\Sigma_{\rm g}\sim 10^{3}\ \rm g\  cm^{-2}$, then $\Sigma_{5}(m_{\bullet}/M_{\odot})\approx 10$. For the prograde cases shown in Figs.~\ref{fig:friction e03 three evo} and \ref{fig:friction e07 three evo}, with the estimated orbital period of $T_{\rm obt}\sim 10^{4}\ M_{\bullet}$, the corresponding number of orbital periods involved in the capture process is $N_{\rm cap}\sim 10^{8}-10^9$. According to the scaling relation in Table.~\ref{tab:power tau p}, the value $N_{\rm cap}\sim 10^3$ reported in Ref.~\cite{spieksma2025gripdiskdraggingcompanion} is in reasonable agreement with the results obtained in this work.\par

%In Ref.~\cite{Wang2024}, the evolution of the orbital parameters $a,\ e$ and $\iota$ was evaluated using $N$-body simulations. For stellar–SMO systems subject to aero-drag force, the eccentricity decreases for all initial inclination angles, leading to orbital circularization, in agreement with our results. In contrast, for sBH–SMO systems subject to dynamical friction, the eccentricity increases for highly inclined aligned configurations as well as for all anti-aligned cases. In both force models and for all initial conditions, the orbital semi-major axis $a$ and inclination angle $\iota$ decrease with time, with comparable and competing evolution timescales.\par

%In Ref.~\cite{Wang2024}, two specific sets of results of $a,\ e,$ and $\iota$ evolution, obtained from $N-$body simulations, were presented, considering an EMRI system consisting of a central SMBH of mass $M_{\bullet}=10^8\ M_{\odot}$ and an SMO of mass $30\ M_{\odot}$. The first case corresponds to a stellar–SMO system with $a_{\rm ini}=10^{-2}\ \mathrm{pc}$ and $\ e_{\rm ini}=0.7$, corresponding to an initial semi-latus rectum $p_{\rm ini}\sim 10^{-3}\ \mathrm{pc}\sim 10^3\ M_{\bullet}$. The eccentricity decreases for all initial inclination angles, leading to orbital circularization, in agreement with our results. For the $\alpha$-disk model with $\alpha=0.01$ considered in Ref.~\cite{Wang2024}, the surface density at the distance of $\sim 10^{3}\ M_{\bullet}$ is $\Sigma_{5}\sim 10^2$. Adopting $m_{*}=30\ M_{\odot}$ and $R_{*}=(m_{*}/M_{\odot})^{\frac{3}{4}}R_{\odot}$~\cite{Wang2024}, the scaling relation in Table.~\ref{tab:power tau p} implies capture timescales that are $5$–$10$ times longer than those reported in Ref.~\cite{Wang2024}.
%The second case corresponds to an sBH–SMO system with initial parameters $a_{\rm ini}=10^{-3}\ \mathrm{pc}$ and $\ e_{\rm ini}=0.7$, yielding an initial semi-latus rectum $p_{\rm ini}\sim 10^{-4}\ \mathrm{pc}\sim 10^{2}\ M_{\bullet}$. In contrast to the first case, the eccentricity increases for highly inclined aligned configurations as well as for all anti-aligned cases. Adopting the $\alpha-$disk model with $\alpha=0.01$ in Ref.~\cite{Wang2024}, the corresponding surface density is $\Sigma_{5}\sim 10^{3}$. For $m_{\bullet}=30\ M_{\odot}$, the scaling relation in Table.~\ref{tab:power tau p} implies capture timescales that are at the order of magnitude as those reported in Ref.~\cite{Wang2024}.\par
% the scaling relations listed in Table~\ref{tab:power tau p} predict capture timescales that are approximately $0.74$ times those obtained in this work. However, the actual results reported in Ref.~\cite{Wang2024} yield significantly shorter capture timescales, corresponding to $\sim 10^{3}$ orbital periods, or $t_{\rm cap}\sim 6\times 10^{7}\ M_{\bullet}-10^{8}\ M_{\bullet}$, whereas our results give $\sim10^{8}-6\times 10^{8}\ M_{\bullet}$. This implies that the capture timescales reported in Ref.~\cite{Wang2024} are shorter than those inferred from the scaling relations of this work, indicating a discrepancy at the level of an order of magnitude. 
%The second case corresponds to an sBH–SMO system with initial parameters $a_{\rm ini}=10^{-3}\ \mathrm{pc}$ and $\ e_{\rm ini}=0.7$, yielding an initial semi-latus rectum $p_{\rm ini}\sim 10^{-4}\ \mathrm{pc}\sim 10^{2}\ M_{\bullet}$. In contrast to the first case, the eccentricity increases for highly inclined aligned configurations as well as for all anti-aligned cases. Adopting the $\alpha-$disk model used in Ref.~\cite{Wang2024}, the corresponding surface density is $\Sigma_{5}\sim 10^{3}$. Using the scaling relations listed in Table~\ref{tab:power tau p}, the capture timescale is therefore expected to be reduced by a factor of $\sim 3\times 10^{-3}$ relative to that obtained in this work. The actual results reported in Ref.~\cite{Wang2024} give capture timescales of $\sim 10^{4}-10^{6}\ \mathrm{yr}$, corresponding to $10^{9}-6\times 10^{10}\ M_{\bullet}$, compared with $t_{\rm cap}\sim 10^{12}-10^{13}\ M_{\bullet}$ found in this work. These values are therefore consistent at the order-of-magnitude level.\par
%then $\Sigma_{5}\sim 10^{3}$ according to the $\alpha-$disk model with $\alpha=0.01$ shown in Ref.~\cite{Wang2024}, the scaling relations listed in Table~\ref{tab:power tau p} predict capture timescales that are reduced by a factor of $\sim 3\times 10^{-3}$ relative to those obtained in this work. The actual results reported in Ref.~\cite{Wang2024} yield capture timescales of $\sim 10^{4}-10^{6}\ \mathrm{yr}$, corresponding to $10^{9}-6\times 10^{10}\ M_{\bullet}$, compared with $t_{\rm cap}\sim 10^{12}-10^{13}\ M_{\bullet}$ found in this work. These values are therefore consistent at the order-of-magnitude level.\par
%Overall, the evolution of the orbital parameters reported in Ref.~\cite{Wang2024} is consistent with that obtained in this work, while the corresponding capture timescales are shorter. Apart from differences in the adopted models, in the calculations of Ref.~\cite{Wang2024}, the disk surface density $\Sigma_{\rm g}$ increases as the orbit of the SMO shrinks, as implied by the disk model, which may further shorten the capture timescale.\par

%In Ref.~\cite{Rowan_2025}, systems with small initial orbital inclination angles were investigated for a central SMBH of mass $M_{\bullet}=10^7\ M_{\odot}$ and an sBH of mass $m_{\bullet}=25\ M_{\odot}$. The alignment timescales estimated under the dynamical friction model were compared with those obtained using the Bondi–Hoyle–Lyttleton (BHL) model with SMO placed at the initial distance of $0.01\ \mathrm{pc}$, corresponding to $\sim 10^{4}\ M_{\bullet}$, with the disk surface density estimated to be $\Sigma_{\rm g}\approx 1.6\times 10^4\ \mathrm{g\ cm^{-2}}$. The maximum orbital inclination angle considered was $\iota=15^{\circ}$, for which the corresponding capture timescale was $t_{\rm cap}\sim 10^{5}\ \mathrm{yr}$. Based on the calculations presented in this work, we estimate that for an initial inclination angle $\iota_{\rm ini}=15^{\circ}$, the capture timescale is of order $t_{\rm cap}\sim 10^{11}\ M_{\bullet}\sim 10^{5}\ \rm yr$. According to the scaling relations listed in Table.~\ref{tab:power tau p}, the resulting timescale is consistent within the same order of magnitude.\par

%In Refs.~\cite{10.1093/mnras/staa3004,10.1093/mnras/stad1295}, the capture of SMOs was investigated under the assumption of circular orbits.
%In these works, the capture timescale was estimated both using a simplified approach with a fixed semi-major axis $a$ and and via numerical integrations that simultaneously evolve $a$ and the orbital inclination angle $\iota$. 
%In Ref.~\cite{10.1093/mnras/staa3004}, aligned configurations with initial inclination angles $\iota_{\rm ini}<90^{\circ}$ were investigated for both stellar and sBH SMOs orbiting a central SMBH of mass $M_{\bullet}=10^{8}\ M_{\odot}$.
%Focusing on interactions with a thin disk, we compare the results obtained from the SG disk model with those presented in this work.
%In particular, we consider the stellar-SMO systems studied in Ref.~\cite{10.1093/mnras/staa3004}, where the SMOs are initially located at radii of $10^2-10^3\ R_{\rm g}$, corresponding to the disk surface density is $\Sigma_{\rm g}\approx 10^5\ \mathrm{g\ cm^{-2}}$.
%In particular, for an M-dwarf SMO with mass $m_{*}=0.5\ M_{\odot}$ and radius $R_{*}=0.4\ R_{\odot}$, the estimated alignment timescale is $\sim 10^5-10^6\ \mathrm{yr}$;
%At the same initial radii of $10^2-10^3\ R_{\rm g}$, a G-type star with mass $m_{*}=1\ M_{\odot}$ and radius $R_{*}=1\ R_{\odot}$ was also considered in Ref.~\cite{10.1093/mnras/staa3004}, for which the alignment timescale was estimated to be $\sim 10^5\ \mathrm{yr}$;
%An O-type star with $m_{*}=50\ M_{\odot}$ and $R_{*}=15\ R_{\odot}$ was found to align on a shorter timescale of $\sim 10^4\ \mathrm{yr}$;
%For a Red-giant SMO with $m_{*}=1.5\ M_{\odot}$ and $R_{*}=100\ R_{\odot}$, the estimated alignment timescale decreases to $\sim 10^1\ \mathrm{yr}$.
%In this work, we obtain capture timescales of $10^8-10^9\ M_{\bullet}$ corresponding to $\sim 10^3-10^4\ \mathrm{yr}$, for stellar–SMO systems. According to the scaling relations listed in Table.~\ref{tab:power tau p}, the capture timescales reported in Ref.~\cite{10.1093/mnras/staa3004} for different stellar types are consistent, at the order-of-magnitude level, with these analytical estimates.
%At high inclinations and large radii, the capture timescales obtained from numerical evolutions are found to deviate significantly from the approximate estimates derived under the assumption $\mathrm{d}a/\mathrm{d}t=0$~\cite{10.1093/mnras/staa3004}, Moreover, a systematic decrease of the capture time with increasing inclination is observed in the numerical results, a trend that is not fully reproduced in the parameter space explored in this work.
%This behavior may be attributed to the more rapid orbital shrinkage in these regimes, together with the steep increase of the SG disk surface density toward smaller radii. This also explains why, by comparison, the numerical results at smaller radii $10^2-10^3\ R_{\rm g}$ are closer to the approximate estimates obtained under the assumption $\mathrm{d}a/\mathrm{d}t=0$, since the disk surface density varies only weakly in this region.\par
%This may be because, in such regions, the orbital shrinkage proceeds more rapidly, and the surface density of the SG disk model exhibits a more pronounced increase toward smaller radii. This also explains why, in comparison, the numerical results in the smaller-radius regions like $10^2-10^3\ R_{\rm g}$ are closer to the approximate results obtained under the assumption of $da/dt=0$, as the disk surface density remains relatively stable there.\par

%For the sBH-SMO cases, the SMO mass in Ref.~\cite{10.1093/mnras/staa3004} was $10\ M_{\odot}$, the alignment timescale can be better fitted with numerical results in Ref.~\cite{10.1093/mnras/staa3004}, while the analytical estimation was larger.
%For the sBH-SMO cases, with an sBH of $m_{\bullet}=10\ M_{\odot}$, the capture timescale inferred from our calculations is in closer agreement with the numerical results reported in Ref.~\cite{10.1093/mnras/staa3004}. %, whereas the corresponding analytical estimates tend to yield systematically longer timescales.
%As a representative example for comparison, we consider the numerical result with $a_{\rm ini}=10^5\ R_{\rm g}$ and $\iota_{\rm ini}=15^{\circ}$. The disk surface density at this distance is $\Sigma_{5}\sim 10^{-2}$. The numerical results reported in Ref.~\cite{10.1093/mnras/staa3004} give the capture timescale 
%According to the scaling relations listed in Table~\ref{tab:power tau p}, this implies a characteristic alignment timescale approximately $27$ times longer than that obtained in this work, i.e., $t_{\rm align}\sim 2.7\times 10^{7}\ \rm yr$. By contrast, the numerical results reported in Ref.~\cite{10.1093/mnras/staa3004} give $t_{\rm align}\sim 4-5\times 10^6\ \mathrm{yr}$, which is closer to the timescale inferred in this work and therefore shorter than the value predicted by the simple scaling estimate. This discrepancy may also be attributed to the fact that, in the actual numerical calculations of Ref.~\cite{10.1093/mnras/staa3004}, the orbit of the small black hole undergoes significant shrinkage while the disk surface density increases toward smaller radii, thereby reducing the alignment timescale.\par

%A more complete range of orbital inclination angles was explored in Ref.~\cite{10.1093/mnras/stad1295}, extending beyond that considered in Ref.~\cite{10.1093/mnras/staa3004}. For sBH-SMO cases with $\iota_{\rm ini}>90^{\circ}$, numerical results showed that the orbital inclination $\iota$ increased once the initial inclination exceeded a critical value. This behavior contrasts with the inclination-damping trend obtained in this work and reported in most previous studies, and may be attributed to the assumption of circular orbits or to the neglect of additional physical effects.%, such as disk backreaction.
%In Ref.~\cite{10.1093/mnras/stad1295}, the full orbital inclination angle range was considered and studied. In the case of sBH–SMO systems initially anti-aligned with the SG disk, the orbital inclination angle $\iota$ increases when it exceeds a critical value. This result is opposite to the inclination-damping behavior found in this work and in most previous studies, and the explanation was the circular-orbit assumption or from the neglect of additional physical effects such as disk back reaction.

%In this work, we investigate the orbital evolution of the SMO in an EMRI system within a Schwarzschild spacetime by employing the forced-geodesic framework and deriving the secular evolution of the key orbital parameters, the semi-major axis $a$, eccentricity $e$, and orbital inclination angle $\iota$, through a two-phase formulation amenable to an adiabatic approximation treatment.
\section{Conclusion}\label{sec:conclusion}
In this work, we study the orbital evolution of EMRI driven by interactions with accretion disks in a relativistic framework. We derive the secular evolution of the key orbital parameters, the semi-major axis $a$, eccentricity $e$, and orbital inclination angle $\iota$, by adopting the double-phase average formulation~\eqref{eq:average integral1}, which is suitable for adiabatic evolution over the proper time $\tau$.
Building on this relativistic framework, we develop a theoretical model for SMO–disk collisions under the thin-disk assumption and analyze the resulting orbital evolution under two representative interaction mechanisms: aero-drag and dynamical friction forces, which capture the essential features of star and sBH crossings.\par
Under both the aero-drag and dynamical friction forces, $a$ and $\iota$ decrease monotonically in both prograde and retrograde configurations, reflecting orbital-energy dissipation and the tendency toward angular-momentum alignment between the orbit and the disk. The eccentricity $e$ decreases monotonically under aero-drag force, leading to a steadily circularizing orbit, whereas under dynamical friction it can be excited at sufficiently large inclinations in some prograde cases and throughout the retrograde cases. In addition, based on the numerical results, we find that the orbit  becomes nearly circularized as the SMO is captured by the disk or as its orbit shrinks.\par
Based on the SG disk model and the scaling relations derived for different interaction mechanisms, together with the alignment times inferred from the representative cases computed in this work, we estimate the capture timescales for wider orbits around SMBHs of various masses and assess the feasibility of capturing SMOs within the timespan of 1 Myr.
We find only a small fraction of sBHs that are initially close to the SMBH and close to the disk
can be captured by the disk within the typical disk lifetime of active galactic nuclei. Therefore, two-body scatterings between sBHs and stars in the nuclear stellar cluster  play an essential role in randomly kicking EMRIs towards the disk and boosting the capture rate \cite{PhysRevD.105.083005}.



We have also derived a power-law dependence of the capture timescale on the initial semi-latus rectum $p$. Specifically, we find $t_{\rm cap}\sim p^{3/2}$ for the aero-drag force model and $t_{\rm cap}\sim p^{-1/2}$ for the dynamical friction model. These scalings are obtained based on the evolution equation of the orbital inclination angle $\iota$ and the disk-crossing process. 
We then compare the orbital evolution and the capture timescales obtained in this work with those reported in several previous studies. We find that the orbital evolution behavior identified here is broadly consistent with most of the existing literature. Moreover, by applying the scaling relations listed in Table~\ref{tab:power tau p} together with more realistic disk surface density models, the estimated capture timescales can be consistent with previous results at the order-of-magnitude level. 

%Some discrepancies nevertheless remain. For instance, the capture timescale reported in Ref.~\cite{Wang2024} is shorter than that estimated using the method adopted in this work. In addition to differences in the adopted evolutionary models, this discrepancy may be attributed to the disk model employed in Ref.~\cite{Wang2024}, in which the disk surface density increases as the SMO orbit shrinks, enhancing the dissipative interaction and further shortening the capture timescale.
%We then compare the  evolution behavior of the orbital parameters, and also the capture timescales obtained in this paper, with several previous work. We find that the evolution behavior in this work is consistent with most of the previous work, and the capture timescale can be consistent at the order of magnitude, using the scaling relations in Table.~\ref{tab:power tau p} and the more realistic disk surface density model. There are also some discrepancies, such like the capture timescale in Ref.~\cite{Wang2024} is shorter than that estimated from the method in this work, apart from differences in the adopted evolution models, it may be explained with that in the calculations of Ref.~\cite{Wang2024}, the disk surface density $\Sigma_{\rm g}$ increases as the orbit of the SMO shrinks, as implied by the disk model, which may further shorten the capture timescale.


%Compared with the parameter trends obtained in Ref.~\cite{Wang2024}, which were obtained from $N$-body simulations in the Newtonian framework, our results exhibit broadly consistent evolutionary behavior for both the stellar–SMO and sBH–SMO cases. Moreover, by comparing the parameter setups and the corresponding capture-timescale estimates derived from scaling relations with the orbital parameter $p$, we find that the capture timescales obtained in this work are consistent with those of Ref.~\cite{Wang2024} within an order of magnitude, albeit with systematic offsets.
%Compared with several other works, like Refs.~\cite{spieksma2025gripdiskdraggingcompanion,Rowan_2025,10.1093/mnras/staa3004,10.1093/mnras/stad1295}, our estimates of the number of capture cycles and the associated timescales exhibit overall agreement, although some discrepancies remain.
%For example, highly inclined orbits have been found to exhibit little to no progression toward alignment in Ref.~\cite{spieksma2025gripdiskdraggingcompanion}. Shorter capture timescales for sBHs than those obtained here have also been reported in Ref.~\cite{10.1093/mnras/staa3004}, along with inclination-increasing behavior in Ref.~\cite{10.1093/mnras/stad1295}. These differences may stem from simplifications adopted in Newtonian-framework treatments of the disk-crossing process, from uncertainties in parameter conversion, and from the fact that a more accurate timescale computation would require extending our framework to incorporate the dependence of the damping coefficient on the disk surface density and scale height.\par
%\subsection{Summary}


\appendix
\section{Basic form of acceleration components}\label{appendix:form of acceleration}
With the basic acceleration form of Eq.~\eqref{eq: force form}, with $u^{\mu}_{\rm ZAMO}$ substituted with $u^{\mu}_{\rm disk}$, the $r,\ \theta,\ \phi$ components of the acceleration can be expressed as
\begin{equation}
\begin{aligned}
    f^{r}(\psi,\ \chi)&=-\gamma_{\rm gas}(\psi,\ \chi)u^{r}\\
    &=-\gamma_{\rm gas}(\psi,\ \chi)\sqrt{\frac{p-6-2e\cos(\psi-\psi_{0})}{p(p-3-e^2)}}e\sin(\psi-\psi_{0}),
\end{aligned}
\label{eq:ar}
\end{equation}
\begin{equation}
\begin{aligned}
    f^{\theta}(\psi,\ \chi)&=-\gamma_{\rm gas}(\psi,\ \chi)u^{\theta}\\
    &=-\gamma_{\rm gas}(\psi,\ \chi)\frac{[1+e\cos(\psi-\psi_{0})]^2}{p\sqrt{p-3-e^2}}\cdot\frac{z_{1}\sin(\chi-\chi_{0})}{\sqrt{1-z_{1}^{2}\cos^{2}(\chi-\chi_{0})}},
\end{aligned}
\label{eq:atheta}
\end{equation}
\begin{equation}
    \begin{aligned}
        f^{\phi}(\psi,\ \chi)&=-\gamma_{\rm gas}(\psi,\ \chi)\left(u^{\phi}+\frac{u^{\phi}_{\rm gas}}{u^{\rm gas}_{t}u^{t}+u^{\rm gas}_{\phi}u^{\phi}}\right)\\
        &=-\gamma_{\rm gas}(\psi,\ \chi)\left\{\frac{[1+e\cos(\psi-\psi_{0})]^{2}}{p\sqrt{p-3-e^2}}\cdot\frac{\sqrt{1-z_{1}^2}}{1-z_{1}^{2}\cos^{2}(\chi-\chi_{0})}\right.\\
        &\left.+\frac{\sqrt{p-3-e^2}}{p\sqrt{1-z_{1}^{2}}-\frac{p\sqrt{(p-2)^2-4e^2}}{\left[1+e\cos(\psi-\psi_{0})\right]^{\frac{3}{2}}}}\right\}.
    \end{aligned}
\label{eq:aphi}
\end{equation}

\section{Coefficients $C_{r,\theta,\phi}^{(I_0)}$}\label{appendix:coefficients}
The detailed expressions of the coefficients for $p,\ e,\ $ and $z_{1}$ in Eq.~\eqref{eq:parameter evolution ode} are
\begin{equation}
\begin{aligned}
     C^{(p)}_{r}&=-\frac{2(p-3-e^2)}{(p-6)^2-4e^2}\sqrt{\frac{e^2p^3[p-2e\cos(\psi-\psi_{0})-6]}{p-3-e^2}}\\
     &\cdot\sin(\psi-\psi_{0}),
\end{aligned}
\end{equation}
\begin{equation}
    \begin{aligned}
    C^{(p)}_{\theta}&=\frac{2p^2 z_{1}\sin(\chi-\chi_{0})\sqrt{p-3-e^2}}{[(p-6)^2-4e^2]\sqrt{1-z_{1}^{2}\cos(\chi-\chi_{0})^{2}}[1+e\cos(\psi-\psi_{0})]^2}\\
    &\cdot[18-9p+p^2-2e(-3+p)\cos(\psi-\psi_{0})-e^2(-6+p)\\
    &\cdot\cos^{2}(\psi-\psi_{0})+2e^{3}\cos^{3}(\psi-\psi_{0})],
\end{aligned}
\end{equation}
\begin{equation}
    \begin{aligned}
        C^{(p)}_{\phi}&=\frac{2p^2\sqrt{(p-3-e^2)(1-z_{1}^{2})}}{[(p-6)^{2}-4e^{2}][1+e\cos(\psi-\psi_{0})]^{2}}\\
        &\cdot[18-9p+p^{2}-2e(-3+p)\cos(\psi-\psi_{0})-e^{2}(-6+p)\\
        &\cdot\cos^{2}(\psi-\psi_{0})+2e^{3}\cos^{3}(\psi-\psi_{0})],
    \end{aligned}
\end{equation}
\begin{equation}
\begin{aligned}
     C^{(e)}_{r}&=\frac{(p-6-2e^2)\sin(\psi-\psi_{0})}{(p-6)^{2}-4e^2}\\
     &\cdot\sqrt{p(p-3-e^{2})[p-6-2e\cos(\psi-\psi_{0})]},
\end{aligned}
\end{equation}
\begin{equation}
    \begin{aligned}
        C^{(e)}_{\theta}&=\frac{pz_{1}\sin(\chi-\chi_{0})\sqrt{p-3-e^{2}}}{[(p-6)^{2}-4e^{2}]\sqrt{1-z_{1}^{2}\cos^{2}(\chi-\chi_{0})}[1+e\cos(\psi-\psi_{0})]^{2}}\\
        &\cdot[e(12+4e^{2}-10p+p^{2})-2(6+2e^{2}-p)(-3+p)\\
        &\cdot\cos(\psi-\psi_{0})-e(6+2e^{2}-p)(-6+p)\cos^{2}(\psi-\psi_{0})\\
        &+2e^{2}(6+2e^{2}-p)\cos^{3}(\psi-\psi_{0})],
    \end{aligned}
\end{equation}
\begin{equation}
    \begin{aligned}
        C^{(e)}_{\phi}&=\frac{p\sqrt{(p-3-e^{2})(1-z_{1}^{2})}}{[(p-6)^{2}-4e^{2}][1+e\cos(\psi-\psi_{0})]^{2}}\\
        &\cdot [e(12+4e^{2}-10p+p^{2})-2(6+2e^{2}-p)(-3+p)\\
        &\cdot\cos(\psi-\psi_{0})-e(6+2e^{2}-p)(-6+p)\cos^{2}(\psi-\psi_{0})\\
        &+2e^{2}(6+2e^{2}-p)\cos^{3}(\psi-\psi_{0})],
    \end{aligned}
\end{equation}
\begin{equation}
    C^{(z_{1})}_{r}=0,
\end{equation}
\begin{equation}
    C^{(z_{1})}_{\theta}=\frac{p(1-z_{1}^{2})\sin(\chi-\chi_{0})\sqrt{p-3-e^{2}}}{\sqrt{1-z_{1}^{2}\cos^{2}(\chi-\chi_{0})}[1+e\cos(\psi-\psi_{0})]^{2}},
    \label{eq:Cz1theta}
\end{equation}
\begin{equation}
    C^{(z_{1})}_{\phi}=-\frac{pz_{1}\sin^{2}(\chi-\chi_{0})\sqrt{(p-3-e^{2})(1-z_{1}^{2})}}{[1+e\cos(\psi-\psi_{0})]^{2}}.
    \label{eq:Cz1phi}
\end{equation}
For $\psi_{0}$ and $\chi_{0}$, with
\begin{equation}
    \begin{aligned}
        \dv{\psi_{r}}{\tau}&=\dv{\psi}{\tau}-\dv{\psi_{0}}{\tau}=\dv{\psi_{r}}{\tau}\Big|_{\rm geodesic}-\dv{\psi_{0}}{\tau}\\
        &=\frac{1}{(\partial r/\partial\psi_{r})}\left(\dv{r}{\tau}-\pdv{r}{p}\dv{p}{\tau}-\pdv{r}{e}\dv{e}{\tau}\right),
    \end{aligned}
\end{equation}
and
\begin{equation}
    \begin{aligned}
        \dv{\psi_{\theta}}{\tau}&=\dv{\chi}{\tau}-\dv{\chi_{0}}{\tau}=\dv{\psi_{\theta}}{\tau}\Big|_{\rm geodesic}-\dv{\chi_{0}}{\tau}\\
        &=\frac{1}{(\partial\theta/\partial\psi_{\theta})}\left(\dv{\theta}{\tau}-\pdv{\theta}{z_{1}}\dv{z_{1}}{\tau}\right),
    \end{aligned}
\end{equation}
then
\begin{equation}
\begin{aligned}
    C^{(\psi_{0})}_{r}&=-\frac{(p-6)\cos(\psi-\psi_{0})+2e}{e[(p-6)^{2}-4e^{2}]}\\
    &\cdot\sqrt{p(p-3-e^{2})[p-6-2e\cos(\psi-\psi_{0})]},
\end{aligned}
\end{equation}
\begin{equation}
    \begin{aligned}
         C^{(\psi_{0})}_{\theta}&=-\frac{pz_{1}\sin(\chi-\chi_{0})\sin(\psi-\psi_{0})\sqrt{p-3-e^{2}}}{e[(p-6)^{2}-4e^{2}]\sqrt{1-z_{1}^{2}\cos^{2}(\chi-\chi_{0})}[1+e\cos(\psi-\psi_{0})]^{2}}\\
        &\cdot\left\{e[4e^{2}-(p-6)^{2}]\cos(\psi-\psi_{0})\right.\\
        &\left.+(-6+p)(6+e^{2}-2p+e^{2}\cos[2(\psi-\psi_{0})])\right\},
    \end{aligned}
\end{equation}
\begin{equation}
    \begin{aligned}
         C^{(\psi_{0})}_{\phi}&=-\frac{p\sqrt{(p-3-e^{2})(1-z_{1}^{2})}\sin(\psi-\psi_{0})}{e[(p-6)^{2}-4e^{2}][1+e\cos(\psi-\psi_{0})]^{2}}\\
        &\cdot\left\{e[4e^{2}-(p-6)^{2}]\cos(\psi-\psi_{0})\right.\\
        &\left.+(-6+p)(6+e^{2}-2p+e^{2}\cos[2(\psi-\psi_{0})])\right\},
    \end{aligned}
\end{equation}
\begin{equation}
    \begin{aligned}
        C^{(\chi_{0})}_{r}=0,
    \end{aligned}
\end{equation}
\begin{equation}
    \begin{aligned}
        C^{(\chi_{0})}_{\theta}=-\frac{p\sqrt{p-3-e^{2}}(1-z_{1}^{2})\cos(\chi-\chi_{0})}{z_{1}\sqrt{1-z_{1}^{2}\cos^{2}(\chi-\chi_{0})}[1+e\cos(\psi-\psi_{0})]^{2}},
    \end{aligned}
\end{equation}
\begin{equation}
    \begin{aligned}
        C^{(\chi_{0})}_{\phi}=\frac{p\sqrt{(p-3-e^{2})(1-z_{1}^{2})}\sin[2(\chi-\chi_{0})]}{2[1+e\cos(\psi-\psi_{0})]^{2}}.
    \end{aligned}
\end{equation}
\section{Adiabatic approximation form in this work}\label{appendix:average Integral2}
In this work, by explicitly accounting for $\psi_{0}$ and $\chi_{0}$, we adopt an equivalent form of Eq.~\eqref{eq:average integral1}, in which the phases $\psi$ and $\chi$ are integrated from $0$ to $2\pi$, as
\begin{equation}
    \expval{f}_{I_{0}}=\frac{\int_{0}^{2\pi}d\chi\int_{0}^{2\pi}d\psi\frac{1}{\omega_{r}(\psi_{r},\ \mathbf{I})}f(\psi_{r},\ \psi_{\theta},\ \mathbf{I})}{2\pi\int_{0}^{2\pi}d\psi\frac{1}{\omega_{r}(\psi_{r},\ \mathbf{I})}},
    \label{eq:average integral2}
\end{equation}
with $\psi_{r}=\psi-\psi_{0}$, $\psi_{\theta}=\chi-\chi_{0}$, $\mathbf{I}=\{p,\ e,\ z_{1}\}$ and $I_{0}=\{p,\ e,\ z_{1},\ \psi_{0},\ \chi_{0}\}$. 
%\section{SG disk profile}
%\label{appendix:pagn}
%The surface density, scale height of the SG disk model generated with \texttt{pagn} code, \texttt{Sirko} module, with different central SMBH masses as $M_{\bullet}=10^6\ M_{\odot},\ 4\times 10^{6}\ M_{\odot},\ 10^7\ M_{\odot},\ 10^{8}\ M_{\odot},\ 4\times 10^{8}\ M_{\odot}$, and the other parameters: the Eddington ratio $l_{e}=0.5$, viscosity parameter $\alpha=0.1$, hydrogen abundance $X=0.7$, and $\alpha-$disk case with $b=0$, are shown in Fig.~\ref{fig:Sigmag Hdisk pagn}.
% \begin{figure}[htbp]
%     \centering
%     \includegraphics[scale=0.55]{pagn_models/Sigma_g_SMBH_mass_contrast.pdf}
%     \caption{The surface density $\Sigma_{\rm g}$ distribution with the radial distance $r$ from the central SMBH with different masses, for SG disk model.}
%     \label{fig:Sigmag pagn}
% \end{figure}
% \begin{figure}[htbp]
%     \centering
%     \includegraphics[scale=0.55]{pagn_models/H_disk_SMBH_mass_contrast.pdf}
%     \caption{The disk scale height $H_{\rm disk}$ distribution with the radial distance $r$ from the central SMBH with different masses, for SG disk model.}
%     \label{fig:Hdisk pagn}
% \end{figure}

\bibliographystyle{apsrev4-1}
\bibliography{reference}
\end{document}